\documentclass[10pt, a4paper]{article}
\usepackage[utf8]{inputenc}
\usepackage[T1]{fontenc}
\usepackage[english]{babel}
\usepackage{amsmath, amssymb, amsfonts, amsthm}
\usepackage{mathrsfs}
\usepackage{graphicx}
\usepackage[hidelinks]{hyperref}
\usepackage{geometry}
\usepackage{mathtools}
\usepackage{braket}
\usepackage{bm}
\usepackage{array}
\usepackage{booktabs}
\usepackage{setspace}
\usepackage{float}

\geometry{a4paper, top=2.5cm, bottom=2.5cm, left=2cm, right=2cm, headheight=14pt}

% --- Updated Commands (Academic Standard) ---
\newcommand{\rrt}{\text{RRT}}  % Referential Relativity Theory
\newcommand{\Tmu}{T_{\mu}}      % Temporal vector field
\newcommand{\Tmunu}{\mathcal{T}_{\mu\nu}} % Temporal force tensor
\newcommand{\Lagrange}{\mathscr{L}}
\newcommand{\BNP}{\text{BNP}}   % Baryonic Neutrality Principle
\newcommand{\Tr}{\text{Tr}}
\newcommand{\Sint}{\mathcal{S}_{\text{int}}}
\newcommand{\Hamil}{\mathcal{H}}
\newcommand{\rhoS}{\rho_S}
\newcommand{\eff}{\text{eff}}
\newcommand{\grad}{\text{grad}}
\newcommand{\Jomega}{J(\omega)}
\newcommand{\Jomegath}{J(\omega, \theta)}
\newcommand{\xiphi}{\xi_\Phi}  % Scalar coupling
\newcommand{\xiT}{\xi_T}       % Vector coupling

\newtheorem{derivation}{Derivation}
\newtheorem{theorem}{Theorem}
\newtheorem{lemma}[theorem]{Lemma}
\newtheorem{corollary}[theorem]{Corollary}

\begin{document}

% --- PAGE 1: TITLE PAGE ---
\begin{titlepage}
    \vspace*{\fill}
    \begin{center}
        {\LARGE \textbf{Referential Relativity Theory - Volume II}} \\
        \vspace{0.6cm}
        {\Large \textbf{Field Formalism and the Dynamics of Open \\ Quantum Systems}} \\
        \vspace{2.5cm}
        
        {\large \textbf{Jean C. Cortez}} \\
        \textit{Independent Researcher and Philosopher} \\
        \vspace{1.2cm}
        
        \textbf{January 2026}
    \end{center}
    \vspace*{\fill}
\end{titlepage}

\newpage

% --- PAGE 2: ABSTRACT ---
\thispagestyle{empty}
\vspace*{\fill}
\begin{abstract}
\setstretch{1.2} 
\noindent This paper presents the complete mathematical derivations of the \textbf{Referential Relativity Theory} (RRT). The formulation is constructed from a covariant action that unifies Einstein-Hilbert geometry with two fundamental temporal fields: the scalar causal flux field $\Phi$ and the anisotropic vector field $T_\mu$, which respectively define cosmic acceleration and temporal directionality. We introduce the physical couplings $\xi_\Phi = \delta = 0.3084$ (resolving the $H_0$ tension) and $\xi_T$ (related to the critical acceleration $a_0 = 1.2001 \times 10^{-10}$ m/s²). 

We derive the total Lagrangian structure, the modified field equations, the Baryonic Neutrality Principle (BNP) in its final form, and the system-reservoir formalism for the vector field $T_\mu$ with a gradient ($D = D_0 \cdot z$) and \textbf{cosmic precession} ($\theta_T(z) = \theta_0 + \omega_p z$). The formulation maintains local compatibility with General Relativity via the BNP, while establishing that the anisotropy—validated through spectroscopic auditing—points to an anti-relation of \textbf{30.36$\sigma$} in the critical stratum of SDSS DR16Q ($z \approx 1.7$), confirming the precession of the causal frame at large scales. Furthermore, we introduce the \textbf{Baryonic Shielding Function} $K(\rho) = 1 - \tanh(\rho/\rho_{crit})$, which governs the phase transition between the stagnation regime (General Relativity) and the viscous regime (RRT). This function reconciles the null results observed in local experiments (LAGEOS-2, CERN) with deep dynamical anomalies. We demonstrate that vacuum viscosity acts as a \textbf{Cortez Refractive Index} $\eta_C(z)$, inducing a phase delay in both light and gravitational waves. This prediction is validated through the auditing of the CASTLES catalog, where RRT provides Theoretical Refractive Convergence without dark matter halos, and predicts a measurable gravitational phase delay for GWTC events.
\end{abstract}

% --- Technical Glossary: Volume II (Field Formalism) ---
\vspace{3cm}
\noindent \textbf{Technical Glossary (RRT Field Dynamics):} \\
\textbf{Baryonic Shielding Function $K(\rho)$:} A sigmoid scaling mechanism that exponentially suppresses temporal viscosity in high-density regions, recovering local Lorentzian invariance. \\
\textbf{Causal Stratigraphy Metric (CSM):} An anisotropic redefinition of the spacetime interval that incorporates a linear phase gradient $D_0 \cdot z$ and precessional trajectories. \\
\textbf{Cortez Refractive Index $\eta_C(z)$:} A redshift-dependent function describing the quantum vacuum as a dissipative medium that induces phase delays in metric and electromagnetic information. \\
\textbf{Causal Reservoir:} The conceptual framework treating the fundamental vector field $T_\mu$ as a non-stationary thermodynamic bath that couples to quantum observables. \\
\textbf{System-Reservoir Formalism:} The mathematical derivation using Feynman-Vernon functionals to describe how "Pure Duration" causes decoherence and dissipation in baryonic matter. \\
\textbf{Causal Stagnation:} The physical state of Phase 1/2 regimes where high gravitational potentials "lock" the temporal tensor, rendering cosmological anisotropy undetectable localy.
\vspace*{\fill}

\newpage

\section*{BLOCK 1: Lagrangian Formulation and Fundamental Couplings of RRT}

\section{Derivation of the Total Referential Relativity Action}

\begin{derivation}[Covariant Action and Field Equations]
The total RRT action synthesizes geometry, matter, and causally active temporal fields, with separate couplings for scalar (acceleration) and vector (anisotropy) effects.

\textbf{Step 1: Total Covariant Action}
\begin{equation}
S_{\text{RRT}} = \int d^4x \sqrt{-g} \left[ \frac{R}{16\pi G} + \Lagrange_\Phi + \Lagrange_T + \Lagrange_{\text{mat}} + \Lagrange_{\text{int}} \right]
\tag{1.1}
\end{equation}

\textbf{Step 2: Lagrangian Components} (Updated)
\begin{itemize}
\item \textbf{Einstein-Hilbert (Geometry)}: $\Lagrange_{\text{EH}} = \dfrac{R}{16\pi G}$
\item \textbf{Scalar Temporal Field (Causal Flux)}:
\begin{equation}
\Lagrange_\Phi = -\frac{1}{2} g^{\mu\nu} \nabla_\mu \Phi \nabla_\nu \Phi - V(\Phi)
\tag{1.2}
\end{equation}
where $V(\Phi)$ is the potential generating the effective cosmic acceleration (replacing dark energy).
\item \textbf{Vector Temporal Field (Fundamental Anisotropy)}:
\begin{equation}
\Lagrange_T = -\frac{1}{4} F_{\mu\nu} F^{\mu\nu} + \frac{m_T^2}{2} T_\mu T^\mu
\tag{1.3}
\end{equation}
with $F_{\mu\nu} = \nabla_\mu T_\nu - \nabla_\nu T_\mu$. The mass $m_T$ may be null or infinitesimal; the critical factor is the existence of a non-zero expectation value $\langle T_\mu \rangle \neq 0$.
\item \textbf{Conventional Matter}: $\Lagrange_{\text{mat}}$
\item \textbf{Causal Interaction and Topological Coupling}:
\begin{equation}
\Lagrange_{\text{int}} = \xiphi \, \Phi \, \mathscr{O}_{\text{mass scale}} + \xiT \, T_\mu J^\mu_{(\text{causal})} + \lambda_{CS} T_\mu K^\mu
\tag{1.4}
\end{equation}
where $K^\mu = \epsilon^{\mu\nu\alpha\beta} (A_\nu F_{\alpha\beta} + \frac{2}{3} A_\nu A_\alpha A_\beta)$ is the \textbf{Chern-Simons Current}. 
\end{itemize}

\textbf{Particle Shielding}: The coupling $\lambda_{CS} T_\mu K^\mu$ identifies the causal vector as the \textbf{Goldstone Boson} of Lorentz symmetry breaking. This term ensures that temporal directionality is communicated to the gauge sector of the Standard Model, allowing RRT to predict cosmic birefringence effects and leptonic asymmetries. This explains why the $T_\mu$ field is suppressed at high energies (LHC) but becomes dominant at cosmological scales where the causal flux density compensates for the topological coupling.

\textbf{Physical Justification}: The separation into $\xiphi$ and $\xiT$ addresses the hierarchy problem and allows for independent calibration from cosmological data. The term $\xiphi \Phi \mathscr{O}_{\text{mass scale}}$ couples the temporal flux density to the mass scale of matter (resolving $H_0$), while $\xiT T_\mu J^\mu$ couples the causal vector to a conserved current associated with matter distribution (including the effective dark component). This prediction underwent stress testing via the BASS DR2 catalog (local AGNs), yielding a result of only \textbf{-0.34$\sigma$}. This data, near statistical nullity at local scales ($z < 0.1$), serves as definitive experimental confirmation of the BNP, proving that baryonic matter in high-density systems is shielded against temporal precession, ensuring the recovery of General Relativity in the nearby universe.

\textbf{Step 3: The Phase Transition Lagrangian}
To formalize the Baryonic Neutrality Principle (BNP), we define the effective Lagrangian density governing local interaction through a sigmoid shielding function:

\begin{equation}
\Lagrange_{\eff} = \Lagrange_{\text{GR}} + K(\rho) \cdot \Lagrange_{\rrt}
\tag{1.4a}
\end{equation}

The \textbf{Shielding Function $K(\rho)$} is phenomenologically defined as:
\begin{equation}
K(\rho) = 1 - \tanh\left( \frac{\rho}{\rho_{crit}} \right)
\tag{1.5}
\end{equation}

This function ensures that:
\begin{itemize}
    \item \textbf{Saturated Regime (Phase 1/2):} In high densities (Solar System, Laboratory), $\rho \gg \rho_{crit} \implies K(\rho) \to 0$. RRT is suppressed, recovering General Relativity.
    \item \textbf{Viscous Regime (Phase 3):} In the intergalactic vacuum, $\rho \ll \rho_{crit} \implies K(\rho) \to 1$. The viscous term $\Lagrange_{\rrt}$ dominates, generating dark matter and dark energy effects.
    \textbf{Note on Effective Density ($\rho_{eff}$):}
For the calculation of the shielding $K(\rho)$, density must be interpreted as the total effective energy density ($T_{00}$), including the local gravitational field energy:
\begin{equation}
\rho_{eff} = \rho_{baryon} + \frac{g^2}{8\pi G c^2}
\end{equation}
In the Solar System and terrestrial orbits (LAGEOS), the gravitational field contribution maintains $\rho_{eff} \gg \rho_{crit}$, ensuring the regime remains in \textbf{Phase 1 (Saturated)}, shielding local experiments.
\end{itemize}

\begin{equation}
G_{\mu\nu} = 8\pi G \left( T_{\mu\nu}^{(\text{mat})} + T_{\mu\nu}^{(\Phi)} + T_{\mu\nu}^{(T)} + T_{\mu\nu}^{(\text{int})} \right)
\tag{1.5}
\end{equation}
\textbf{The Causal Stratigraphy Metric (CSM)}: 
RRT requires a transition from the isotropic FLRW metric to the CSM, where the line element incorporates the vector anisotropy $T_\mu$. The interval $ds^2$ is redefined to capture the linear phase gradient:
\begin{equation}
ds^2 = -(1 + \Psi) dt^2 + a(t)^2 \exp\left[ D_0 \cdot z \cdot \cos(\theta - \theta_T(z)) \right] \gamma_{ij} dx^i dx^j
\tag{1.5.1}
\end{equation}

From this metric, we derive the \textbf{Cortez Refractive Index} $\eta_C$:
\begin{equation}
\eta_C(z) = 1 + \beta \cdot \ln(1 + z)
\tag{1.5.2}
\end{equation}
Physically, $\eta_C$ describes the vacuum as a refractive medium. This effect explains the time delay observed in gravitational lensing without the need for dark matter halos, providing a solid theoretical foundation for the mass-geometry mismatch in critical systems such as PG1115+080.

The exponential term acts as a \textbf{Causal Refraction Factor}. This implies that the luminosity distance ($d_L$) is no longer a purely scalar function of redshift but becomes a directional tensor. This modification shields RRT against the Hubble Tension, as it demonstrates that the discrepancy between SH0ES and Planck is a frame effect within the causal flux.

\textbf{Step 4: Derived Energy-Momentum Tensors}
\begin{align}
T_{\mu\nu}^{(\Phi)} &= \nabla_\mu \Phi \nabla_\nu \Phi - g_{\mu\nu} \left[ \frac{1}{2} \nabla_\alpha \Phi \nabla^\alpha \Phi + V(\Phi) \right]
\tag{1.6} \\
T_{\mu\nu}^{(T)} &= F_{\mu\alpha} F^\alpha_{\ \nu} - \frac{1}{4} g_{\mu\nu} F_{\alpha\beta} F^{\alpha\beta} + m_T^2 \left( T_\mu T_\nu - \frac{1}{2} g_{\mu\nu} T_\alpha T^\alpha \right)
\tag{1.7} \\
T_{\mu\nu}^{(\text{int})} &= \xiphi \Phi \, T_{\mu\nu}^{(\text{mat, scale})} + \frac{\xiT}{2} \left( T_\mu J_\nu + T_\nu J_\mu - g_{\mu\nu} T_\alpha J^\alpha \right)
\tag{1.8}
\end{align}

\textbf{Step 5: Causal Conservation and Modified Bianchi Identities}
Unlike General Relativity, where $\nabla^\mu G_{\mu\nu} = 0$ in isolation, RRT imposes an extended conservation. We define the Einstein-Cortez Tensor as $\mathcal{G}_{\mu\nu} = G_{\mu\nu} + \kappa \mathcal{T}_{\mu\nu}$. Theoretical consistency requires that the total system energy (geometry + temporal field) be conserved:
\begin{equation}
\nabla^\mu (G_{\mu\nu} + \kappa \mathcal{T}_{\mu\nu}) = 0
\tag{1.9}
\end{equation}
This formulation ensures that any energy flux from the temporal field is locally compensated by metric curvature. This proves that the "extra" energy detected in galaxies is not created ex nihilo, but is a momentum transfer from the temporal tensorial field to baryonic matter via the $\xi$ coupling, satisfying the generalized Bianchi Identities.
\end{derivation}

\textbf{Note on Stability and Ghost Modes}: To ensure the theoretical shielding of the vector sector, we require the $\xi_T$ coupling to satisfy the stability condition $c_s^2 = 1 - \frac{2\xi_T^2}{3} > 0$. With the calibrated value of $\xi_T \approx 10^{-3}$, the propagation of perturbations in the temporal field $T_\mu$ is strictly subluminal and free of negative energies, ensuring that RRT does not violate quantum unitarity.

% ==============================================================================
% NEW SECTION: MATHEMATICAL PROOF OF PHASE CONTINUITY
% ==============================================================================
\sloppy {Continuous Phase Transition Dynamics: The Landau-Cortez Derivation}

To demonstrate that RRT and General Relativity are not distinct theories artificially joined, but rather asymptotic regimes of a single geometric functionality, we derive here the system evolution from a Unified Variational Principle.

We define the \textbf{Order Parameter $\psi(x)$}, which quantifies the "fluidity" of time:
\begin{itemize}
    \item $\psi = 0$: Solid/Rigid State (Lorentzian Isotropy).
    \item $\psi = 1$: Fluid/Viscous State (Causal Anisotropy).
\end{itemize}

\begin{figure}[H]
    \centering
    \includegraphics[width=0.85\textwidth]{grafico37a.png}
    \caption{\textbf{Vacuum Viscosity Ladder:} Visual representation of RRT phase regimes. As density ($\rho$) decreases, the order parameter $\psi$ transitions from 0 (rigid vacuum/Einstein) to 1 (viscous fluid vacuum/Cortez), revealing frame dynamics at cosmological scales.}
    \label{fig:viscosity_ladder}
\end{figure}

\subsection{The Master Action Functional}
The total action of the universe is controlled by a Lagrangian density that blends the regimes through a phase potential $V(\psi, \rho)$:

\begin{equation}
S_{Total} = \int d^4x \sqrt{-g} \left[ \underbrace{(1 - \psi) \frac{R}{16\pi G}}_{\text{Rigid Geometry (GR)}} + \underbrace{\psi \mathcal{L}_{RRT}}_{\text{Fluid Geometry}} + \mathcal{L}_{mat} - V(\psi, \rho) \right]
\end{equation}

Where the Transition Potential $V(\psi, \rho)$ follows the Landau-Ginzburg form, controlled by the local energy density $\rho$:
\begin{equation}
V(\psi, \rho) = \frac{\lambda}{2} (\rho - \rho_{crit}) \psi^2 + \frac{\beta}{4} \psi^4
\end{equation}

The nature of the universe at any point is determined by minimizing this action: $\frac{\delta S}{\delta \psi} = 0$.

\subsection{Step 1: Derivation of Phase 1 (High-Density Saturation)}
Consider the interior of a stellar system or a particle collider (CERN), where $\rho \gg \rho_{crit}$.
The coefficient $(\rho - \rho_{crit})$ becomes a large positive number. To minimize the potential energy $V$, the system forces the order parameter to collapse:
\begin{equation}
\frac{\partial V}{\partial \psi} \approx \lambda \rho \psi = 0 \implies \psi \to 0
\end{equation}

Substituting $\psi = 0$ into the Master Action:
\begin{equation}
S_{Phase 1} = \int d^4x \sqrt{-g} \left[ \frac{R}{16\pi G} + \mathcal{L}_{mat} \right]
\end{equation}
\textbf{Result:} We recover exactly the Einstein-Hilbert Action. Temporal viscosity vanishes, gravity becomes purely metric, and the Standard Model is shielded. This mathematically explains the null results observed in LAGEOS and CMS.

\subsection{Step 2: Derivation of Phase 2 (Galactic Transition)}
Consider the galactic halo region, where the baryon density falls to $\rho \approx \rho_{crit}$.
The quadratic term of the potential vanishes $(\rho - \rho_{crit} \approx 0)$. The system enters symmetry breaking, allowing $\psi$ to fluctuate ($\psi > 0$).
The equation of motion for the gravitational potential $\Phi_N$ receives a mixed contribution:
\begin{equation}
\nabla^2 \Phi_{Total} = 4\pi G \rho + \nabla \cdot (\psi \vec{\mathcal{T}})
\end{equation}
Where the second term acts as a "support pressure." Solving for virial equilibrium:
\begin{equation}
v_{orbital}^4 = G M a_0 \quad (\text{Tully-Fisher Relation})
\end{equation}
\textbf{Result:} The flat rotation curve emerges naturally from the mixture of phases, without the need for Dark Matter. The "extra mass" is, in fact, the momentum of the temporal fluid ($\psi \vec{\mathcal{T}}$) coupling to matter.

\subsection{Step 3: Derivation of Phase 3 (The Viscous Vacuum)}
In deep intergalactic space (Voids), $\rho \ll \rho_{crit}$. The potential coefficient becomes negative, inverting the energy parabola. The new energy minimum shifts to $\psi \to 1$ (Viscous Saturation).
The Action is now dominated by the fluid term:
\begin{equation}
S_{Phase 3} \approx \int d^4x \sqrt{-g} \left[ \mathcal{L}_{RRT} - V_{min} \right]
\end{equation}
The equation of state for this temporal fluid generates a stress tensor with effective negative pressure:
\begin{equation}
P_{eff} = -\rho_{vac} c^2
\end{equation}
\textbf{Result:} This negative pressure drives the accelerated expansion of the universe. Dark Energy is not an inexplicable cosmological constant $\Lambda$, but the thermodynamic manifestation of time flowing freely ($\psi=1$) in the vacuum.

\subsection{Derivation Conclusion}
We have demonstrated that Phases 0, 1, 2, and 3 are continuous solutions of the same Master Action. Singularity is avoided because, as extreme densities are approached, the variable $\psi$ automatically adjusts to dissipate or absorb degrees of freedom, keeping the system finite and conservative.


\section{The Recovery of General Relativity as a Limiting Case}

RRT does not invalidate General Relativity (GR); it encompasses it as the regime of \textbf{Causal Stagnation}. The Einstein-Hilbert action is identically recovered when the divergence of the causal vector tends to zero.

\begin{equation}
\lim_{\nabla_\mu T^\mu \to 0} \Lagrange_{RRT} = \frac{R}{16\pi G} + \Lagrange_{mat}
\end{equation}

% --- BNP SHIELDING GRAPH INSERTION ---
\begin{figure}[H]
    \centering
    \includegraphics[width=0.95\textwidth]{grafico19a.png}
    \caption{\textbf{Baryonic Neutrality Principle (BNP) Attenuation Curve.} The graph illustrates the phase transition governed by the local density $\rho_b$. In the deep vacuum regime (left, red), causal coupling is maximal ($\chi \approx 1$), allowing the detection of 30.36$\sigma$ anisotropies in SDSS. In the high-density regime (right, blue), such as the Solar System, coupling tends to zero ($\chi \to 0$), recovering the precision of General Relativity and explaining the null results observed in the LAGEOS satellite.}
    \label{fig:pnb_shielding}
\end{figure}
% --------------------------------------------

Physically, baryonic matter acts as a thermo-causal insulator. In regions of high gravitational potential ($\Phi > \Phi_{crit}$), the function $\chi(\Phi) \to 0$ "deactivates" the tensor $\mathcal{T}_{\mu\nu}$ for the local observer. This explains why inertial reference satellites like LAGEOS-2 report statistical nullity (\textbf{0\% field activation}), while galactic edge dynamics (where $\Phi$ is low) reveal the full drag of the temporal field.

\subsection{Preservation of Nucleosynthesis (BBN)}
RRT maintains the precision of the standard model in the primordial universe due to thermal saturation. At energies above 1 GeV, the particle collision rate exceeds the precession frequency $\omega_p$, forcing the isotropization of the vacuum. This ensures that light element abundances remain unaltered, while the $\mathcal{T}_{\mu\nu}$ field emerges as the dominant force only in the low-density regime (Phase 3).

\section{EFT Regime Map and the Order Parameter in RRT}

\begin{derivation}[UV→IR Transition Governed by $\nabla_\mu T^\mu$]
\textbf{Introduction}: The phase transition in RRT is characterized by the divergence of the causal vector $\nabla_\mu T^\mu$, which acts as an order parameter. This section describes the complete Effective Field Theory (EFT) regime map.

\textbf{Step 1: Definition of the Divergence Operator}
\begin{equation}
\mathcal{O}_{\text{div}}(x) = \nabla_\mu T^\mu(x)
\tag{1.10}
\end{equation}

\textbf{Step 2: UV Regime (High Energy - $E \gg \Lambda_{\text{RRT}}$)}
In the ultraviolet regime, $\nabla_\mu T^\mu$ is a relevant operator:
\begin{equation}
\Delta_{\text{UV}}(\nabla_\mu T^\mu) = 3 - \epsilon, \quad \epsilon > 0
\tag{1.11}
\end{equation}

In this regime:
\begin{itemize}
\item Lorentz symmetry breaking is explicit and strong.
\item The causal vector $T_\mu$ defines absolute preferred directions.
\item The anomalous dimension $\epsilon$ is determined by the renormalization group.
\item This regime corresponds to the first fractions of a second after the Big Bang.
\end{itemize}

\textbf{Step 3: Intermediate Regime (Energy $\sim \Lambda_{\text{RRT}}$)}
At the characteristic scale $\Lambda_{\text{RRT}} \sim H_0$:
\begin{equation}
\Delta_{\text{int}}(\nabla_\mu T^\mu) = 4 \quad \text{(marginal operator)}
\tag{1.12}
\end{equation}
In this critical regime:
\begin{itemize}
\item Causal coupling becomes marginal.
\item The Law of Linear Causal Coupling $D = D_0 \cdot z$ emerges.
\item The observational gradient $D_0$ is calibrated.
\item Cosmic precession mechanism $\omega_p(z)$ appears.
\end{itemize}

\textbf{Step 4: IR Regime (Low Energy/Local - $E \ll \Lambda_{\text{RRT}}$)}
In the infrared regime:
\begin{equation}
\Delta_{\text{IR}}(\nabla_\mu T^\mu) = 4 + \delta, \quad \delta > 0
\tag{1.13}
\end{equation}
In this limit:
\begin{itemize}
\item $\nabla_\mu T^\mu \to 0$ is an irrelevant operator.
\item General Relativity is recovered locally.
\item The Baryonic Neutrality Principle (BNP) becomes exact.
\item Solar System gravity tests are satisfied.
\end{itemize}

\textbf{Step 5: Renormalization Group Equation}
\begin{equation}
\mu \frac{d}{d\mu} \nabla_\mu T^\mu = \beta(g) \nabla_\mu T^\mu
\tag{1.14}
\end{equation}
with beta function:
\begin{equation}
\beta(g) = \mu \frac{dg}{d\mu} = - \epsilon g + \frac{g^3}{16\pi^2} + \mathcal{O}(g^5)
\tag{1.15}
\end{equation}

\textbf{Step 6: Fixed Point and Phase Transition}
The fixed point occurs when:
\begin{align}
\beta(g_*) &= 0
\tag{1.16} \\
\frac{d\beta}{dg}\bigg|_{g_*} &> 0 \quad \text{(attractive fixed point in the IR)}
\tag{1.17}
\end{align}
The solution:
\begin{equation}
g_*^2 = 16\pi^2 \epsilon
\tag{1.18}
\end{equation}

\textbf{Step 7: Connection with Cosmological Observations}
In the modified FRW metric:
\begin{align}
\nabla_\mu T^\mu &= \dot{T}_0 + 3H T_0 + \frac{\partial T_i}{\partial x^i}
\tag{1.19} \\
&= D_0 \cdot H(z) \cdot (1+z) \cdot \cos(\theta - \theta_T(z))
\tag{1.20}
\end{align}
This relationship directly links the theoretical order parameter ($\nabla_\mu T^\mu$) to observations ($D_0$, $\omega_p$).

\textbf{Step 8: Physical Interpretation}
\begin{itemize}
\item \textbf{UV Regime}: The early universe possessed strong temporal directionality.
\item \textbf{Intermediate Regime}: The current era, where we measure gradient and precession.
\item \textbf{IR Regime}: The cosmic future, where anisotropy becomes irrelevant.
\end{itemize}

\textbf{Step 9: Complete Phase Map}
\begin{table}[ht]
\centering
\caption{RRT Regime Map}
\begin{tabular}{lccc}
\toprule
\textbf{Regime} & \textbf{$\Delta(\nabla_\mu T^\mu)$} & \textbf{Behavior} & \textbf{Cosmic Epoch} \\
\midrule
UV (High Energy) & $< 4$ (relevant) & Strong Lorentz breaking & $t < 10^{-32}$ s \\
Critical & $= 4$ (marginal) & Linear coupling & Current era ($z \sim 0-10$) \\
IR (Low Energy) & $> 4$ (irrelevant) & GR recovery & Local scale ($z \approx 0$) \\
\bottomrule
\end{tabular}
\end{table}

\textbf{Step 10: Mathematical Shielding}
The UV→IR transition is not a regularization artifact but a fundamental characteristic:
\begin{itemize}
\item In the UV: $\nabla_\mu T^\mu$ is a relevant operator (Symmetry Breaking).
\item Intermediate: Marginal operator (Law of Linear Causal Coupling).
\item In the IR: Irrelevant operator, where $\nabla_\mu T^\mu \to 0$.
\end{itemize}
This structure protects the theory against criticisms of inconsistency between regimes.
\end{derivation}

\section{The Cortez Metric and the Unified Mass Equation}

The immediate physical consequence of tensorial anisotropy is the redefinition of observed mass ($M_{obs}$). We demonstrate that the effective gravitational potential is not purely Einsteinian but is amplified by the temporal flux density in the observer's reference frame.

\begin{derivation}[The Unified Cortez Equation]
The calculated mass ($M_{calc}$) that an observer detects via gravitational lensing or stellar dynamics, without the need to invoke dark matter, is given by:
\begin{equation}
    M_{calc} = M_{\star} \cdot \left( 1 + [(D_0 \cdot z)^2 \cdot e^{z/\pi}] + \Omega_p \right)
\end{equation}
Where:
\begin{itemize}
    \item $M_{\star}$ is the actual detected baryonic mass (stars and gas).
    \item $D_0 = 0.794$ is the linear causal coupling constant.
    \item $e^{z/\pi}$ represents the accumulated temporal density factor (dominates at high $z$).
    \item $\Omega_p$ is the local vorticity factor, responsible for rotation anomalies in nearby galaxies.
\end{itemize}
This equation replaces the dark matter density term ($\rho_{DM}$) with a geometric frame adjustment, reducing residual error by orders of magnitude relative to the $\Lambda$CDM model.
\end{derivation}

\section{Derivation of the Baryonic Neutrality Principle as a State of Minimum Causal Entropy}

\begin{derivation}[BNP as the "State of Minimum Causal Entropy Production" in the Feynman-Vernon Functional]
\textbf{Introduction}: The Baryonic Neutrality Principle (BNP) is not an ad hoc imposition but emerges naturally as the stationary state of minimum causal entropy production within the open systems formalism.

\textbf{Step 1: Feynman-Vernon Influence Functional for RRT}
\begin{equation}
\mathcal{F}[A_S, A_S'] = \Tr_T \left[ e^{-\frac{i}{\hbar} \mathcal{S}_{\text{int}}[A_S]} \rho_T e^{\frac{i}{\hbar} \mathcal{S}_{\text{int}}[A_S']} \right]
\tag{3.1}
\end{equation}

\textbf{Step 2: Causal Entropy Production}
We define the causal entropy production per unit time:
\begin{equation}
\dot{\Sigma}_{\text{causal}} = -k_B \frac{d}{dt} D_{\text{KL}}(\rho_S(t) || \rho_{\text{BNP}})
\tag{3.2}
\end{equation}
where $D_{\text{KL}}$ is the Kullback-Leibler divergence and $\rho_{\text{BNP}}$ is the state satisfying the BNP.

\textbf{Step 3: Minimum Condition}
The BNP corresponds to the state that minimizes causal entropy production:
\begin{equation}
\frac{\delta \dot{\Sigma}_{\text{causal}}}{\delta \rho_S} \bigg|_{\rho_S = \rho_{\text{BNP}}} = 0
\tag{3.3}
\end{equation}

\textbf{Step 4: Variational Derivation}
Consider the effective action:
\begin{equation}
\Gamma[\rho_S] = \int dt \left[ \dot{\Sigma}_{\text{causal}}(t) + \lambda(t)(\Tr\rho_S(t) - 1) \right]
\tag{3.4}
\end{equation}
Varying with respect to $\rho_S$:
\begin{align}
\frac{\delta \Gamma}{\delta \rho_S} &= \frac{\partial \dot{\Sigma}_{\text{causal}}}{\partial \rho_S} + \lambda(t) I = 0
\tag{3.5} \\
\Rightarrow \rho_{\text{BNP}} &= \arg\min_{\rho_S} \dot{\Sigma}_{\text{causal}}(\rho_S)
\tag{3.6}
\end{align}

\textbf{Step 5: Explicit Solution in Lindblad Formalism}
For the master equation:
\begin{equation}
\frac{d\rho_S}{dt} = -i[H_S, \rho_S] + \mathcal{L}_{\text{diss}}(\rho_S)
\tag{3.7}
\end{equation}
The BNP state satisfies:
\begin{equation}
\mathcal{L}_{\text{diss}}(\rho_{\text{BNP}}) = 0 \quad \text{and} \quad [H_S, \rho_{\text{BNP}}] = 0
\tag{3.8}
\end{equation}

\textbf{Step 6: General Form of the BNP State}
\begin{equation}
\rho_{\text{BNP}} = \frac{e^{-\beta_{\text{eff}} H_S}}{Z} \otimes P_{\text{neutral}}
\tag{3.9}
\end{equation}
where $P_{\text{neutral}}$ projects onto the zero causal charge subspace.

\textbf{Step 7: Baryonic Neutrality Condition}
For baryonic matter:
\begin{equation}
\Tr[\rho_{\text{BNP}} \, \xiT T_\mu J^\mu_{(\text{bar})}] = 0
\tag{3.10}
\end{equation}
This condition emerges as the Euler-Lagrange equation of the entropy functional.

\textbf{Step 8: Thermodynamic Interpretation}
The BNP is the open-systems analogue of the state of thermodynamic equilibrium:
\begin{itemize}
\item \textbf{Thermal Equilibrium}: Minimizes thermal entropy production.
\item \textbf{BNP}: Minimizes causal entropy production.
\end{itemize}

\textbf{Step 9: Stability and Attractivity}
The BNP state is stable and attractive:
\begin{equation}
\frac{d}{dt} D_{\text{KL}}(\rho_S(t) || \rho_{\text{BNP}}) \leq 0
\tag{3.11}
\end{equation}
with equality only at $\rho_S = \rho_{\text{BNP}}$.

\textbf{Step 10: Connection to the Regime Map}
\begin{itemize}
\item In the UV: Causal entropy production is maximal.
\item In the IR: $\dot{\Sigma}_{\text{causal}} \to 0$, BNP becomes exact.
\item Intermediate: Smooth transition between regimes.
\end{itemize}

\textbf{Conclusion}: The BNP emerges inevitably as the stationary state of matter immersed in the causal reservoir, being the state of minimum causal entropy production within the Feynman-Vernon influence functional. This derivation shields the theory against criticisms of BNP artificiality.
\end{derivation}

\section{Methodological Update: Hubble Detrending and Clean Gradient}

\begin{derivation}[Methodological Transition for Causal Signal Isolation]
\textbf{Step 1: The Methodological Turning Point}
The major change in the 2025 analysis was the transition from raw magnitude analysis to \textbf{Hubble Detrending (Expansion Cleanup)}. Previously, the anisotropy signal was "contaminated" by the expansion of the universe; now, we isolate the pure vacuum gradient.

\textbf{Step 2: Previous Approach (Obsolete)}
\begin{equation}
\Delta m_{\text{old}} = Mag_{\text{obs}} - \langle Mag \rangle
\tag{3.12}
\end{equation}
This method captured the acceleration of the universe as if it were anisotropy, contaminating the causal signal.

\textbf{Step 3: New Approach (RRT 2025) - Hubble Detrending}
\begin{equation}
\text{Residual} = Mag_{\text{obs}} - (5 \log_{10} z)
\tag{3.13}
\end{equation}
This removes the logarithmic Hubble expansion curve. What remains is only the directional oscillation caused by the Causal Vector $T^\mu$.

\textbf{Step 4: Blinding Test Implementation}
To ensure the signal is not statistical noise, we implemented:
\begin{enumerate}
\item Code executed on \textbf{random universes} (Monte Carlo simulations).
\item Verification that the sigma found is not compatible with random fluctuations.
\item Cross-validation with multiple independent catalogs.
\end{enumerate}

\textbf{Step 5: Hubble Detrending Result}
The observed gradient $D = D_0 \cdot z$ emerges clearly after removing Hubble expansion, confirming that:
\begin{itemize}
\item The signal is not an artifact of $H_0$ calibration.
\item Anisotropy is intrinsic to the cosmic vacuum.
\item The Law of Linear Causal Coupling ($D = D_0 \cdot z$) is empirically validated.
\end{itemize}

\textbf{Step 6: Cosmic Stratigraphy}
This method allows for **cosmic stratigraphy**: redshift layers reveal the temporal evolution of causal coupling, showing that the deeper we look into the vacuum, the more evident the coupling becomes.
\end{derivation}

\section{The Golden Triad: Updated Catalog Results}

\begin{derivation}[Final Calibration of RRT Parameters]
\textbf{Step 1: Corrected Numbers Dashboard (The Golden Triad)}
These are the official figures included in all RRT 2025 manuscripts:

% --- CORRECTED BLOCK FOR VOLUME II ---
\begin{table}[ht] 
\centering
\caption{RRT Audit Results 2025 (The Golden Triad)}
\label{tab:golden_triad}
\resizebox{\textwidth}{!}{%
\begin{tabular}{lcccc}
\toprule
\textbf{Catalog} & \textbf{Objects (N)} & \textbf{Significance} & \textbf{Result} & \textbf{Verdict} \\
\midrule
Planck (CMB) & -- & \textbf{98.36\%} & Ecliptic Geometry & Primordial Unification \\ 
SDSS DR16Q & 218,210 & \textbf{30.36$\sigma$} & Parity Inversion & Causal Resonance \\ 
SPARC (Galax.) & 175 & \textbf{1.33\%} & Mean Residual Error & End of Dark Matter \\
BASS DR2 & 94 & \textbf{-0.34$\sigma$} & Local Isotropy & Phase 1 Confirmation \\
\bottomrule
\end{tabular}}
\end{table}


\begin{figure}[H]
    \centering
    \includegraphics[width=0.8\textwidth]{grafico14a.png}
    \caption{Geometry-Mass Correlation in SLACS. The oscillation of mass mismatch as a function of the Cortez Angle proves that Dark Matter is a pseudo-force resulting from the non-metricity of time.}
    \label{fig:slacs_v2}
\end{figure}
% --- END OF CORRECTED BLOCK ---

\textbf{Step 2: Technical Note on Deep Surveys and Decoherence}
In deep surveys such as DES, the exact vector rotation cancels the global linear signal, generating the expected phase decoherence. This validates the RRT prediction that the causal reference frame precesses and accumulates phase shifts with redshift.

\textbf{Step 3: Inter-Catalog Consistency}
\begin{itemize}
\item \textbf{Pantheon+} (1,426 SNe): 23.24$\sigma$ - the gold standard for fine calibration and definitive confirmation of the causal gradient.
\end{itemize}

\textbf{Step 4: Implications for $\xiphi$ and $\xiT$ Couplings}
These updated values lead to:
\begin{align}
\xiT &= (2.18 \pm 0.03) \times 10^{-3} \quad \text{(recalibrated with new $D_0$)}
\tag{3.14} \\
a_0^{\text{RRT}} &= 1.2001 \times 10^{-10} \, \text{m/s}^2 \quad \text{(0.09\% precision)}
\tag{3.15}
\end{align}

\textbf{Step 5: Validation of the Law of Linear Causal Coupling}
The linear gradient $D = D_0 \cdot z$ is confirmed with unprecedented precision:
\begin{equation}
\chi^2_{\text{red}} = 1.02 \quad \text{(ideal linear fit)}
\tag{3.16}
\end{equation}
rejecting quadratic or exponential models at $>5\sigma$.

\textbf{Step 6: Updated Cosmic Precession}
With new data, the precession is refined:
\begin{equation}
\theta_T(z) = \theta_0 + \omega_p z \quad \text{with} \quad \omega_p \text{ as the Cosmic Precession Rate}
\tag{3.17}
\end{equation}
\end{derivation}

\section{Derivation of the Modified Field Equations}

\begin{derivation}[Full Action Variation]
\textbf{Step 1: Equation for the Scalar Field $\Phi$}
Varying $S_{\text{RRT}}$ with respect to $\Phi$:
\begin{equation}
\nabla_\mu \nabla^\mu \Phi - \frac{dV}{d\Phi} = - \xiphi \mathscr{O}_{\text{mass scale}}
\tag{4.1}
\end{equation}
In a weak-field and homogeneous approximation:
\begin{equation}
\ddot{\Phi} + 3H\dot{\Phi} + \frac{dV}{d\Phi} = \xiphi \rho_m
\tag{4.2}
\end{equation}

\textbf{Step 2: Equation for the Vector Field $T_\mu$}
Varying with respect to $T_\mu$:
\begin{equation}
\nabla_\nu F^{\nu\mu} + m_T^2 T^\mu = - \xiT J^\mu_{(\text{causal})}
\tag{4.3}
\end{equation}
In a homogeneous but anisotropic cosmological background, with $\langle T_\mu \rangle = (T_0(t), \mathbf{T}(t))$:
\begin{equation}
\ddot{\mathbf{T}} + 3H\dot{\mathbf{T}} + m_T^2 \mathbf{T} = - \xiT \mathbf{J}_{(\text{causal})}
\tag{4.4}
\end{equation}

\textbf{Step 3: Modified Friedmann Equations}
From the (0,0) and (i,i) components of (1.5):
\begin{align}
\left(\frac{\dot{a}}{a}\right)^2 &= \frac{8\pi G}{3} \left( \rho_m + \rho_\Phi + \rho_T \right) - \frac{k}{a^2}
\tag{4.5} \\
\frac{\ddot{a}}{a} &= -\frac{4\pi G}{3} \left( \rho_m + \rho_\Phi + \rho_T + 3P_\Phi + 3P_T \right)
\tag{4.6}
\end{align}
with:
\begin{align}
\rho_\Phi &= \frac{1}{2}\dot{\Phi}^2 + V(\Phi), \quad P_\Phi = \frac{1}{2}\dot{\Phi}^2 - V(\Phi)
\tag{4.7} \\
\rho_T &= \frac{1}{2}\dot{T}_i\dot{T}^i + \frac{1}{4}F_{ij}F^{ij} + \frac{m_T^2}{2}(T_0^2 + T_i T^i)
\tag{4.8} \\
P_T &= \frac{1}{2}\dot{T}_i\dot{T}^i + \frac{1}{6}F_{ij}F^{ij} - \frac{m_T^2}{6}(T_0^2 + T_i T^i)
\tag{4.9}
\end{align}

\textbf{Step 4: Baryonic Neutrality Principle (BNP) - Final Formulation}
For ordinary baryonic matter, we impose:
\begin{equation}
\int d^3 x \sqrt{-g} \, J^0_{(\text{causal})} \bigg|_{\text{bar.}} = 0 \quad \Rightarrow \quad \langle \xiT T_\mu J^\mu_{(\text{bar})} \rangle_{\text{local}} \approx 0
\tag{4.10}
\end{equation}
This ensures that Solar System gravity tests (PPN limits) are satisfied, while allowing for cosmologically significant couplings for effective dark matter and fields.

\textbf{Step 5: Causal Screening and Stagnation Mechanism} In regions of stable high baryonic density, the energy density of matter acts as a "stagnation anode," preventing local vacuum precession. This maintains $\chi(\Phi) \approx 0$, inducing a state of causal stagnation (*stagnation*) that shields the $\mathcal{T}_{\mu\nu}$ tensor, making anisotropy undetectable in satellites such as LAGEOS-2 (\textbf{0\% causal drag}). The function $\chi(\Phi)$ is not a free parameter but a description of the inertial phase transition: activation occurs only in the low-acceleration regime ($a < a_0$), where vacuum viscosity resumes its metric dominance over matter, revealing time precession as the true engine of missing mass in galaxies.
\end{derivation}

\section{Derivation of the Temporal Force Tensor \texorpdfstring{$\mathcal{T}_{\mu\nu}$}{T\_munu} in RRT}

\begin{derivation}[From the Variational Formalism]
In RRT, the temporal force tensor $\Tmunu$ is not an independent fundamental field but emerges as an effective combination of vector and interaction terms.

\textbf{Step 1: Variational Definition}
\begin{equation}
\mathcal{T}_{\mu\nu} = \xi_T (\bar{T}_ \mu J_\nu + \bar{T}_v J_\mu) + \lambda_{CS} \mathcal{K}_{\mu\nu}
\tag{5.2}
\end{equation}
where $\mathcal{K}_{\mu\nu}$ is the tensor derived from the Chern-Simons current. The inclusion of this topological term shields the particle sector, identifying $T_\mu$ as the Goldstone Boson of Causality.

\textbf{Step 2: Explicit Form for Small Fluctuations}
For small deviations from the isotropic background, with $\Tmu = \bar{T}_\mu + \delta T_\mu$ and $J^\mu = \bar{J}^\mu + \delta J^\mu$:
\begin{equation}
\Tmunu(x) \approx \xiT \left( \bar{T}_\mu \delta J_\nu(x) + \bar{T}_\nu \delta J_\mu(x) \right) + \xiT \left( \delta T_\mu(x) \bar{J}_\nu + \delta T_\nu(x) \bar{J}_\mu \right)
\tag{5.2}
\end{equation}

\textbf{Step 3: Coupled Conservation Law}
\begin{equation}
\nabla^\mu \Tmunu = - \xiT F_{\nu\alpha} J^\alpha_{(\text{causal})} \approx - \xiT m_T^2 \bar{T}_\nu \delta \rho
\tag{5.3}
\end{equation}

\textbf{Step 4: Physical Interpretation}
$\Tmunu$ is the tensorial manifestation of the "force of time" that mediates causal interaction between:
\begin{itemize}
\item The preferred direction of time ($\bar{T}_\mu$).
\item Matter/energy fluctuations ($\delta J_\nu$).
\item Fluctuations of the vector field itself ($\delta T_\mu$).
\end{itemize}
This tensor is responsible for the system-reservoir coupling in the quantum formulation.
\end{derivation}

\section{Derivation of the RRT Hamiltonian Density}

\begin{derivation}[Complete Legendre Transform]
\textbf{Step 1: Updated Conjugate Momenta}
\begin{align}
\pi_\Phi &= \frac{\partial \Lagrange}{\partial \dot{\Phi}} = \dot{\Phi}
\tag{6.1} \\
\pi_{T^i} &= \frac{\partial \Lagrange}{\partial \dot{T}_i} = -F^{0i} = -\dot{T}^i + \nabla^i T^0
\tag{6.2} \\
\pi_{T^0} &= \frac{\partial \Lagrange}{\partial \dot{T}_0} \approx 0 \quad \text{(primary constraint for $m_T=0$)}
\tag{6.3}
\end{align}

\textbf{Step 2: Total Hamiltonian Density}
\begin{equation}
\Hamil_{\text{RRT}} = \Hamil_{\text{EH}} + \Hamil_\Phi + \Hamil_T + \Hamil_{\text{int}}
\tag{6.4}
\end{equation}
with:
\begin{align}
\Hamil_\Phi &= \frac{1}{2} \pi_\Phi^2 + \frac{1}{2} (\nabla \Phi)^2 + V(\Phi)
\tag{6.5} \\
\Hamil_T &= \frac{1}{2} (\pi_{T^i})^2 + \frac{1}{4} F_{ij} F^{ij} + \frac{m_T^2}{2} (T_0^2 + T_i T^i) + \text{gauge terms}
\tag{6.6} \\
\Hamil_{\text{int}} &= -\xiphi \Phi \mathscr{O}_{\text{mass scale}} - \xiT T_\mu J^\mu_{(\text{causal})}
\tag{6.7}
\end{align}

\textbf{Step 3: Weak-Field Approximation for $\Hamil_{\text{int}}$}
When $|\xiT T_\mu J^\mu| \ll |\Hamil_T|$ and $|\xiphi \Phi \mathscr{O}| \ll |\Hamil_\Phi|$:
\begin{equation}
\Hamil_{\text{int}} \approx -\xiT T_\mu J^\mu_{(\text{causal})} \quad \text{(dominant for decoherence effects)}
\tag{6.8}
\end{equation}

\textbf{Step 4: Operational Form for Quantum Calculations}
For a local quantum system with operator $A_S$ coupling to the field:
\begin{equation}
H_{\text{int}} = -\xiT \int_{\mathcal{V}} d^3x \, A_S(x) T_\mu(x) J^\mu(x)
\tag{6.9}
\end{equation}
where $\mathcal{V}$ is the volume of the system.
\end{derivation}

\section{Derivation of the Preliminary Quantization of Temporal Fields}

\begin{derivation}[Canonical Quantization in Flat Spacetime]
\textbf{Step 1: Promotion to Operators}
\begin{align}
\hat{\Phi}(\mathbf{x}) &\rightarrow \text{scalar operator}
\tag{7.1} \\
\hat{T}_\mu(\mathbf{x}) &\rightarrow \text{vector operator}
\tag{7.2}
\end{align}

\textbf{Step 2: Commutation Relations}
For the vector field (Proca or Stueckelberg-type quantization):
\begin{align}
[\hat{T}_i(\mathbf{x}), \hat{\pi}_{T^j}(\mathbf{y})] &= i\hbar \delta_{ij} \delta^3(\mathbf{x} - \mathbf{y})
\tag{7.3} \\
[\hat{T}_0(\mathbf{x}), \hat{\pi}_{T^0}(\mathbf{y})] &= i\hbar \delta^3(\mathbf{x} - \mathbf{y}) \quad \text{(if $m_T \neq 0$)}
\tag{7.4}
\end{align}

\textbf{Step 3: Causal Vacuum State}
The vacuum state $|0_T\rangle$ is defined by:
\begin{equation}
\hat{T}_\mu |0_T\rangle = \bar{T}_\mu |0_T\rangle
\tag{7.5}
\end{equation}
where $\bar{T}_\mu \neq 0$ is the non-zero expectation value that breaks Lorentzian isotropy. This is a spontaneous symmetry breaking.

\textbf{Step 4: Mode Expansion with Non-Zero Background}
We write:
\begin{equation}
\hat{T}_\mu(x) = \bar{T}_\mu + \delta\hat{T}_\mu(x)
\tag{7.6}
\end{equation}
where $\delta\hat{T}_\mu$ are fluctuation operators with the usual expansion:
\begin{equation}
\delta\hat{T}_\mu(x) = \sum_{\lambda=1}^3 \int \frac{d^3k}{(2\pi)^3} \frac{1}{\sqrt{2\omega_k}} \left[ \epsilon_\mu^{(\lambda)}(k) a^{(\lambda)}(\mathbf{k}) e^{ik\cdot x} + \text{h.c.} \right]
\tag{7.7}
\end{equation}
with $[a^{(\lambda)}(\mathbf{k}), a^{(\lambda')\dagger}(\mathbf{k}')] = \delta_{\lambda\lambda'} \delta^3(\mathbf{k}-\mathbf{k}')$.

\textbf{Step 5: Physical Interpretation}
The background $\bar{T}_\mu$ defines the preferred direction of time (causal axis). The fluctuations $\delta\hat{T}_\mu$ are the quanta of the temporal field that mediate causal interaction with quantum systems.

\textbf{Step 6: Propagation Speed and GW170817 Compatibility}
In the high-energy regime (UV), the effective mass of the $T_\mu$ field becomes negligible ($m_T \to 0$), ensuring that the propagation speed of causal perturbations coincides with the speed of light ($c_{causal} = c$).

The wave equation for coupled perturbations takes the form:
\begin{equation}
\square \delta T_\mu \approx 0
\tag{7.8}
\end{equation}

Consequently, RRT is compatible with observations of the GW170817 event (binary neutron star merger), where the arrival time difference between gravitational waves and gamma rays was less than 2 seconds over a travel distance of 130 million years. RRT predicts that temporal "drag" affects the wave's phase, but not its group velocity in vacuum.
\end{derivation}

\section{Derivation of the System-Reservoir Form with Gradient}

\begin{derivation}[Reduction to the Canonical Formalism with Temporal Evolution]
This is the central derivation connecting RRT with the open systems formalism, incorporating the gradient $D = D_0 \cdot z$.

\textbf{Step 1: Causal Reservoir Operator}
We define the reservoir operator as the component of the vector field that couples linearly:
\begin{equation}
\hat{B}_{\text{AC}}(x) = \xiT \, \delta\hat{T}_\mu(x) \bar{J}^\mu
\tag{8.1}
\end{equation}
Or, more generally, if $A_S$ is the system operator:
\begin{equation}
\hat{B}_{\text{AC}}(x) = \xiT \, \delta\hat{T}_\mu(x) \quad \text{(for $A_S \otimes \hat{B}_{\text{AC}}$ coupling)}
\tag{8.2}
\end{equation}

\textbf{Step 2: Modal Expansion with Redshift Dependence}
Due to cosmic expansion, modes have time-dependent frequencies. In conformal coordinates $\eta$:
\begin{equation}
\delta\hat{T}_\mu(\eta, \mathbf{x}) = \sum_{\lambda} \int \frac{d^3k}{(2\pi)^3} \frac{1}{\sqrt{2\omega_k(\eta)}} \left[ \epsilon_\mu^{(\lambda)}(k) a^{(\lambda)}(\mathbf{k}) e^{i\mathbf{k}\cdot\mathbf{x} - i\int^{\eta} \omega_k(\eta') d\eta'} + \text{h.c.} \right]
\tag{8.3}
\end{equation}

\textbf{Step 3: Interaction Hamiltonian in System-Reservoir Form}
For a system with operator $A_S$:
\begin{align}
H_{\text{int}}(t) &= A_S \otimes \sum_{\lambda} \int \frac{d^3k}{(2\pi)^3} \left[ g_k^{(\lambda)}(t) a^{(\lambda)}(\mathbf{k}) e^{-i\int^t \omega_k(t') dt'} + \text{h.c.} \right]
\tag{8.4} \\
g_k^{(\lambda)}(t) &= \xiT \frac{\epsilon_\mu^{(\lambda)}(k) \bar{J}^\mu}{\sqrt{2\omega_k(t)}} \tilde{A}_S(\mathbf{k})
\tag{8.5}
\end{align}
where $\tilde{A}_S(\mathbf{k})$ is the spatial Fourier transform of the system operator.

\textbf{Step 4: Incorporation of the $D(z)$ Gradient}
The observed gradient $D = D_0 \cdot z$ emerges naturally when considering the temporal evolution of the coupling coefficients:
\begin{equation}
g_k^{(\lambda)}(t) = g_k^{(\lambda)}(0) \cdot (1 + z(t))^{\gamma_{\text{grad}}}
\tag{8.6}
\end{equation}
with $\gamma_{\text{grad}} \approx 1$ to reproduce $D = D_0 \cdot z$. The linear dependence on $z$ arises from integration along the line of sight:
\begin{equation}
\Delta \mu_T(z, \theta) = \int_0^z \frac{dz'}{H(z')} \, \xiT(z') \, \bar{T}_\mu(z') k^\mu \cdot \cos\theta \approx D_0 \cdot z \cdot \cos\theta
\tag{8.7}
\end{equation}
where $k^\mu$ is the photon four-vector.

\textbf{Note:} The value of $D_0$ is now calibrated by the updated Golden Triad data (Table \ref{tab:golden_triad}), with a weighted average $D_0 = 0.794 \pm 0.003$.

\textbf{Step 5: Cosmic Precession}
The direction of the causal vector evolves with redshift:
\begin{equation}
\hat{T}_\mu(z) = R_{\mu\nu}(z) \hat{T}_\nu(0)
\tag{8.8}
\end{equation}
with the rotation matrix $R_{\mu\nu}(z)$ generating the observed precession $\theta_T(z) = \theta_0 + \omega_p z$.

\textbf{Geometric Implication of Precession}: 
Since the direction of the causal axis $\theta_T(z)$ rotates following its precessional rate, the Particle Horizon in RRT assumes the topology of a \textbf{Causal Spiral}. Photons originating from high redshifts do not follow straight Euclidean geodesics but rather precessional trajectories. At $z \approx 1100$ (CMB), the precession frequency exceeds the diffusion scale, resulting in the phase cancellation that preserves the small-scale isotropy measured by the Planck satellite, while allowing for the emergence of the 'Axis of Evil' at large scales ($l < 20$).

\textbf{Step 6: Final Discretized Form}
In finite volume $V$:
\begin{equation}
H_{\text{int}}(t) = \hbar A_S \otimes \sum_{k,\lambda} \left[ g_{k}^{(\lambda)}(t) b_{k}^{(\lambda)} e^{-i\omega_k t} + \text{h.c.} \right]
\tag{8.9}
\end{equation}
with $[b_{k}^{(\lambda)}, b_{k'}^{(\lambda')\dagger}] = \delta_{kk'} \delta_{\lambda\lambda'}$ and coefficients $g_{k}^{(\lambda)}(t)$ carrying gradient and precession information.
\end{derivation}

\textbf{Step 7: Gravitational Phase Delay ($\Delta \tau$)}
Unlike the $\Lambda$CDM model, RRT predicts that $\beta$ viscosity induces an accumulated delay in the propagation of metric information:
\begin{equation}
\Delta \tau = \frac{\beta}{c} \int_0^{D_L} \eta_C(z') dr'
\tag{7.9}
\end{equation}
Theoretical analysis of the gravitational wave propagation confirms that gravity, like light, is subject to the phase impedance of the viscous vacuum, predicting a systematic bias in deep-space measurements.

\section{Derivation of the Baryonic Neutrality Principle (BNP) in the Quantum Formulation}

\begin{derivation}[Zero Coupling Condition for Baryonic Matter]
\textbf{Step 1: Operational Formulation}
For any state $|\psi_{\text{bar}}\rangle$ of pure baryonic matter:
\begin{equation}
\langle \psi_{\text{bar}} | \xiT \hat{T}_\mu \hat{J}^\mu_{(\text{causal})} | \psi_{\text{bar}} \rangle = 0
\tag{9.1}
\end{equation}

\textbf{Step 2: Implementation via Symmetry}
This can be ensured by imposing that the baryonic current operator be orthogonal to the causal vector:
\begin{equation}
\hat{J}^\mu_{(\text{bar})} \bar{T}_\mu = 0
\tag{9.2}
\end{equation}
Or alternatively, that the total temporal charge be zero:
\begin{equation}
Q_T^{\text{(bar)}} = \int d^3x \sqrt{-g} \, \hat{J}^0_{(\text{bar})} = 0
\tag{9.3}
\end{equation}

\textbf{Step 3: Consequences for the S-Matrix}
The scattering matrix between baryonic states via the exchange of $T_\mu$ field quanta is suppressed:
\begin{equation}
\langle f_{\text{bar}} | S | i_{\text{bar}} \rangle_{T-\text{exchange}} \approx 0
\tag{9.4}
\end{equation}

\textbf{Step 4: Preservation in Quantum Loops}
The BNP must be preserved by renormalization, which requires counter-terms to respect condition (9.1). This can be ensured by an appropriate discrete or gauge symmetry.

\textbf{Step 5: Physical Interpretation}
The BNP is the mechanism allowing the temporal field to have strong cosmological effects (dark matter/energy) while remaining "invisible" in laboratory gravity tests.
\end{derivation}

% --- BLOCK 2: Influence Functional and Correlation Functions in RRT ---

\section*{BLOCK 2: Influence Functional and Correlation Functions in RRT}

\textbf{Conceptual Transition}: We now move from the fundamental formulation of RRT to the derivation of the influence functional, which encodes the effects of the causal reservoir (the $T_\mu$ field with gradient and precession) on the quantum system's dynamics. The primary novelty is the explicit dependence on redshift/cosmic time in the coupling coefficients.

\section{Derivation of the Temporal Interaction Action with Cosmic Evolution}

\begin{derivation}[Temporal Integration of the Hamiltonian with Gradient]
The interaction action now incorporates the temporal evolution of coefficients due to cosmic expansion and the gradient.

\textbf{Step 1: Definition of the Interaction Action}
\begin{equation}
\mathcal{S}_{\text{int}} = \int_{t_0}^{t} d\tau \, H_{\text{int}}(\tau)
\tag{10.1}
\end{equation}

\textbf{Step 2: Substitution of the Updated System-Reservoir Form}
Using (8.9) with time-dependent coefficients:
\begin{align}
\mathcal{S}_{\text{int}} &= \int_{t_0}^{t} d\tau \, \hbar A_S(\tau) \otimes \sum_{k,\lambda} \left[ g_{k}^{(\lambda)}(\tau) b_{k}^{(\lambda)} e^{-i\int^{\tau} \omega_k(\tau') d\tau'} + \text{h.c.} \right]
\tag{10.2} \\
&= \hbar \sum_{k,\lambda} \int_{t_0}^{t} d\tau \, A_S(\tau) \left[ g_{k}^{(\lambda)}(\tau) b_{k}^{(\lambda)} e^{-i\Omega_k(\tau)} + g_{k}^{(\lambda)*}(\tau) b_{k}^{(\lambda)\dagger} e^{i\Omega_k(\tau)} \right]
\tag{10.3}
\end{align}
where $\Omega_k(\tau) = \int^{\tau} \omega_k(\tau') d\tau'$ is the accumulated phase.

\textbf{Step 3: Parametric Form of the Coefficient with Gradient}
Based on (8.6), we assume:
\begin{equation}
g_{k}^{(\lambda)}(t) = g_{k0}^{(\lambda)} \cdot (1 + z(t))^{\gamma_k}
\tag{10.4}
\end{equation}
To reproduce $D = D_0 \cdot z$, we require $\gamma_k \approx 1$ for low-frequency modes dominating large-scale interaction. More precisely:
\begin{equation}
g_{k}^{(\lambda)}(t) = g_{k0}^{(\lambda)} \left[ 1 + \frac{D_0}{c} \int_{t_0}^{t} H(t') dt' \cdot \cos\theta_k \right]
\tag{10.5}
\end{equation}
where $\theta_k$ is the angle between $\mathbf{k}$ and the causal vector direction $\hat{T}$.

\textbf{Note:} The value of $D_0$ in this expression is calibrated by the Golden Triad ($D_0 = 0.794 \pm 0.003$).

\textbf{Step 4: Operator Separation}
System and reservoir operators commute at equal times:
\begin{equation}
[A_S(\tau), b_{k}^{(\lambda)}] = 0, \quad [A_S(\tau), b_{k}^{(\lambda)\dagger}] = 0
\tag{10.6}
\end{equation}

\textbf{Step 5: Final Form of the Action}
\begin{equation}
\mathcal{S}_{\text{int}} = \hbar \sum_{k,\lambda} \left[ b_{k}^{(\lambda)} I_{k}^{(\lambda)}(t,t_0) + b_{k}^{(\lambda)\dagger} I_{k}^{(\lambda)*}(t,t_0) \right]
\tag{10.7}
\end{equation}
with modified temporal integrals:
\begin{equation}
I_{k}^{(\lambda)}(t,t_0) = \int_{t_0}^{t} d\tau \, A_S(\tau) g_{k}^{(\lambda)}(\tau) e^{-i\Omega_k(\tau)}
\tag{10.8}
\end{equation}

\textbf{Physical Interpretation}: The action represents the accumulated coupling between the system and the causal reservoir, with increasing weight due to the $D_0 \cdot z$ gradient. The $(1+z)^{\gamma_k}$ term encapsulates the "causal memory" accumulated over cosmic history.
\end{derivation}

\sloppy{Derivation of the Influence Functional Definition for Cosmic Reservoir}

\begin{derivation}[Trace over $T_\mu$ Field Degrees of Freedom]
\textbf{Step 1: Updated Formal Definition}
\begin{equation}
\mathcal{F}[A_S, A_S'] = \Tr_{T} \left[ \exp\left( -\frac{i}{\hbar} \mathcal{S}_{\text{int}}[A_S] \right) \rho_T \exp\left( \frac{i}{\hbar} \mathcal{S}_{\text{int}}[A_S'] \right) \right]
\tag{11.1}
\end{equation}
where $\rho_T$ is the state of the $T_\mu$ field (generally a thermal state with temperature $T_{\text{causal}}$).

\textbf{Step 2: Cosmic Reservoir State}
For the $T_\mu$ vector field with non-zero background:
\begin{equation}
\rho_T = \frac{e^{-\beta H_T}}{\Tr[e^{-\beta H_T}]}, \quad \beta = (k_B T_{\text{causal}})^{-1}
\tag{11.2}
\end{equation}
with $H_T$ given by (6.6). The thermal state is defined relative to fluctuations $\delta T_\mu$, not including the background $\bar{T}_\mu$.

\textbf{Step 3: Modified Evolution Operators}
\begin{align}
U_{A_S}(t,t_0) &= \exp\left( -\frac{i}{\hbar} \mathcal{S}_{\text{int}}[A_S] \right)
\tag{11.3} \\
U_{A_S'}(t,t_0) &= \exp\left( \frac{i}{\hbar} \mathcal{S}_{\text{int}}[A_S'] \right)
\tag{11.4}
\end{align}

\textbf{Step 4: Interpretation as Amplitude}
$\mathcal{F}[A_S, A_S']$ represents the transition amplitude of the $T_\mu$ field under the influence of histories $A_S(\tau)$ and $A_S'(\tau)$. The redshift dependence in $g_k(\tau)$ causes the functional to "sense" the cosmic history of the coupling.

\textbf{Step 5: Normalization Condition}
For $A_S = A_S'$, we have:
\begin{equation}
\mathcal{F}[A_S, A_S] = \Tr_T[\rho_T] = 1
\tag{11.5}
\end{equation}

\textbf{Step 6: Conjugation Property}
\begin{equation}
\mathcal{F}[A_S, A_S']^* = \mathcal{F}[A_S', A_S]
\tag{11.6}
\end{equation}

\textbf{Step 7: Connection with the Path Integral in RRT}
In the path integral formulation of RRT, the functional arises when tracing over $\delta T_\mu$:
\begin{equation}
K(q_f, t; q_i, 0) = \int \mathscr{D}q \, e^{\frac{i}{\hbar} S_S[q]} \mathcal{F}[q] \cdot e^{\frac{i}{\hbar} S_{\text{back}}[\bar{T}]}
\tag{11.7}
\end{equation}
where $S_{\text{back}}[\bar{T}]$ is the action of the $\bar{T}_\mu$ background, providing additional direction-dependent phase factors (precession effect).
\end{derivation}

\section{Derivation of the Second-Order Expansion with Time-Dependent Coefficients}

\begin{derivation}[Development of the Dyson Series for Non-Stationary Coupling]
We expand the exponentials to second order, now with time-dependent $g_k(t)$ coefficients.

\textbf{Step 1: Expansion of Evolution Operators}
\begin{align}
e^{-\frac{i}{\hbar} \mathcal{S}_{\text{int}}[A_S]} &\approx 1 - \frac{i}{\hbar} \mathcal{S}_{\text{int}}[A_S] - \frac{1}{2\hbar^2} \mathcal{S}_{\text{int}}[A_S]^2
\tag{12.1} \\
e^{\frac{i}{\hbar} \mathcal{S}_{\text{int}}[A_S']} &\approx 1 + \frac{i}{\hbar} \mathcal{S}_{\text{int}}[A_S'] - \frac{1}{2\hbar^2} \mathcal{S}_{\text{int}}[A_S']^2
\tag{12.2}
\end{align}

\textbf{Step 2: Operator Product}
\begin{align}
&e^{-\frac{i}{\hbar} \mathcal{S}_{\text{int}}[A_S]} \rho_T e^{\frac{i}{\hbar} \mathcal{S}_{\text{int}}[A_S']} \\
&= \rho_T + \frac{i}{\hbar} \mathcal{S}_{\text{int}}[A_S'] \rho_T - \frac{i}{\hbar} \rho_T \mathcal{S}_{\text{int}}[A_S] \\
&\quad - \frac{1}{2\hbar^2} \mathcal{S}_{\text{int}}[A_S']^2 \rho_T + \frac{1}{\hbar^2} \mathcal{S}_{\text{int}}[A_S'] \rho_T \mathcal{S}_{\text{int}}[A_S] - \frac{1}{2\hbar^2} \rho_T \mathcal{S}_{\text{int}}[A_S]^2 \\
&\quad + \mathcal{O}(\xiT^3)
\tag{12.3}
\end{align}

\textbf{Step 3: Trace over the Reservoir}
Taking the trace and using $\Tr_T[b_{k}^{(\lambda)} \rho_T] = 0$:
\begin{align}
\mathcal{F}[A_S, A_S'] &\approx 1 - \frac{1}{2\hbar^2} \Tr_T[\mathcal{S}_{\text{int}}[A_S]^2 \rho_T] \\
&\quad + \frac{1}{\hbar^2} \Tr_T[\mathcal{S}_{\text{int}}[A_S] \rho_T \mathcal{S}_{\text{int}}[A_S']] \\
&\quad - \frac{1}{2\hbar^2} \Tr_T[\rho_T \mathcal{S}_{\text{int}}[A_S']^2]
\tag{12.4}
\end{align}

\textbf{Step 4: Term Simplification}
Defining expectation values $\langle \cdot \rangle_T = \Tr_T[\cdot \rho_T]$:
\begin{align}
\Tr_T[\mathcal{S}_{\text{int}}[A_S]^2 \rho_T] &= \langle \mathcal{S}_{\text{int}}[A_S]^2 \rangle_T
\tag{12.5} \\
\Tr_T[\rho_T \mathcal{S}_{\text{int}}[A_S']^2] &= \langle \mathcal{S}_{\text{int}}[A_S']^2 \rangle_T
\tag{12.6} \\
\Tr_T[\mathcal{S}_{\text{int}}[A_S] \rho_T \mathcal{S}_{\text{int}}[A_S']] &= \langle \mathcal{S}_{\text{int}}[A_S] \mathcal{S}_{\text{int}}[A_S'] \rangle_T
\tag{12.7}
\end{align}

\textbf{Step 5: Compact Form}
\begin{equation}
\mathcal{F}[A_S, A_S'] \approx 1 - \frac{1}{2\hbar^2} \left( \langle \mathcal{S}_{\text{int}}[A_S]^2 \rangle_T + \langle \mathcal{S}_{\text{int}}[A_S']^2 \rangle_T - 2\langle \mathcal{S}_{\text{int}}[A_S] \mathcal{S}_{\text{int}}[A_S'] \rangle_T \right)
\tag{12.8}
\end{equation}
\end{derivation}

\section{Derivation of the Expression in Terms of Anisotropic Correlation Functions}

\begin{derivation}[Identification of Correlation Functions with Gradient]
\textbf{Step 1: Expression of the Action in Terms of Operators}
Using (10.3):
\begin{align}
\mathcal{S}_{\text{int}}[A_S] &= \int_{t_0}^{t} d\tau \, A_S(\tau) B(\tau)
\tag{13.1} \\
\mathcal{S}_{\text{int}}[A_S'] &= \int_{t_0}^{t} d\tau \, A_S'(\tau) B(\tau)
\tag{13.2}
\end{align}
where we define the reservoir operator:
\begin{equation}
B(\tau) = \hbar \sum_{k,\lambda} \left[ g_{k}^{(\lambda)}(\tau) b_{k}^{(\lambda)} e^{-i\Omega_k(\tau)} + \text{h.c.} \right]
\tag{13.3}
\end{equation}

\textbf{Step 2: Product of Actions}
\begin{align}
\mathcal{S}_{\text{int}}[A_S] \mathcal{S}_{\text{int}}[A_S'] &= \int_{t_0}^{t} d\tau \int_{t_0}^{t} d\tau' \, A_S(\tau) A_S'(\tau') B(\tau) B(\tau')
\tag{13.4}
\end{align}

\textbf{Step 3: Thermal Expectation Values}
\begin{align}
\langle \mathcal{S}_{\text{int}}[A_S]^2 \rangle_T &= \int_{t_0}^{t} d\tau \int_{t_0}^{t} d\tau' \, A_S(\tau) A_S(\tau') \langle B(\tau) B(\tau') \rangle_T
\tag{13.5} \\
\langle \mathcal{S}_{\text{int}}[A_S']^2 \rangle_T &= \int_{t_0}^{t} d\tau \int_{t_0}^{t} d\tau' \, A_S'(\tau) A_S'(\tau') \langle B(\tau) B(\tau') \rangle_T
\tag{13.6} \\
\langle \mathcal{S}_{\text{int}}[A_S] \mathcal{S}_{\text{int}}[A_S'] \rangle_T &= \int_{t_0}^{t} d\tau \int_{t_0}^{t} d\tau' \, A_S(\tau) A_S'(\tau') \langle B(\tau) B(\tau') \rangle_T
\tag{13.7}
\end{align}

\textbf{Step 4: Definition of the Time-Dependent Correlation Function}
\begin{equation}
\mathcal{C}(\tau, \tau') = \langle B(\tau) B(\tau') \rangle_T
\tag{13.8}
\end{equation}
\textbf{Important}: Due to the time dependence in $g_k(\tau)$, $\mathcal{C}(\tau, \tau')$ is no longer a function only of $\tau-\tau'$, but of both times separately. However, we can write:
\begin{equation}
\mathcal{C}(\tau, \tau') = f(\tau, \tau') \, C_0(\tau - \tau')
\tag{13.9}
\end{equation}
where $C_0(\Delta\tau)$ is the stationary correlation function and $f(\tau, \tau')$ captures the evolution due to the gradient:
\begin{equation}
f(\tau, \tau') = \left[ (1+z(\tau))(1+z(\tau')) \right]^{\gamma/2}
\tag{13.10}
\end{equation}
with $\gamma \approx 1$.

\textbf{Step 5: Final Influence Functional Expression}
\begin{align}
\mathcal{F}[A_S, A_S'] &\approx 1 - \frac{1}{2} \int_{t_0}^{t} d\tau \int_{t_0}^{t} d\tau' \, f(\tau, \tau') \big[ A_S(\tau) A_S(\tau') C_0(\tau - \tau') \\
&\quad + A_S'(\tau) A_S'(\tau') C_0^*(\tau - \tau') - 2 A_S(\tau) A_S'(\tau') C_0(\tau - \tau') \big]
\tag{13.11}
\end{align}
\end{derivation}

\section{Derivation of the Exponential Form via Generalized Cumulant Theorem}

\begin{derivation}[Re-exponentiation for Non-Stationary Gaussian Processes]
\textbf{Step 1: Cumulant Theorem for Non-Stationary Gaussian Processes}
For a Gaussian variable $X(t)$ with $\langle X(t) \rangle = 0$ but non-stationary correlations $\langle X(t)X(t') \rangle$:
\begin{equation}
\langle e^{i\int dt \, X(t)h(t)} \rangle = \exp\left[ -\frac{1}{2} \int dt \int dt' \, h(t)h(t') \langle X(t)X(t') \rangle \right]
\tag{14.1}
\end{equation}

\textbf{Step 2: Application to the Influence Functional}
Defining:
\begin{equation}
X = \frac{i}{\hbar} (\mathcal{S}_{\text{int}}[A_S] - \mathcal{S}_{\text{int}}[A_S'])
\tag{14.2}
\end{equation}
We have:
\begin{equation}
\mathcal{F}[A_S, A_S'] = \langle e^{X} \rangle_T = e^{\frac{1}{2}\langle X^2 \rangle_T}
\tag{14.3}
\end{equation}
(Note: for Gaussian variables with zero mean, all cumulants of order higher than 2 are zero).

\textbf{Step 3: Calculation of the Second Cumulant}
\begin{align}
\langle X^2 \rangle_T &= -\frac{1}{\hbar^2} \langle (\mathcal{S}_{\text{int}}[A_S] - \mathcal{S}_{\text{int}}[A_S'])^2 \rangle_T \\
&= -\frac{1}{\hbar^2} \left( \langle \mathcal{S}_{\text{int}}[A_S]^2 \rangle_T + \langle \mathcal{S}_{\text{int}}[A_S']^2 \rangle_T - 2\langle \mathcal{S}_{\text{int}}[A_S] \mathcal{S}_{\text{int}}[A_S'] \rangle_T \right)
\tag{14.4}
\end{align}

\textbf{Step 4: Substitution of Non-Stationary Correlations}
Using (13.5)-(13.7):
\begin{align}
\langle X^2 \rangle_T &= -\frac{1}{\hbar^2} \int_{t_0}^{t} d\tau \int_{t_0}^{t} d\tau' \, f(\tau, \tau') \big[ A_S(\tau) A_S(\tau') C_0(\Delta\tau) \\
&\quad + A_S'(\tau) A_S'(\tau') C_0^*(\Delta\tau) - 2 A_S(\tau) A_S'(\tau') C_0(\Delta\tau) \big]
\tag{14.5}
\end{align}
where $\Delta\tau = \tau - \tau'$.

\textbf{Step 5: Final Exponential Form}
\begin{align}
\mathcal{F}[A_S, A_S'] &\approx 1 - \frac{1}{2} \int_{t_0}^{t} d\tau \int_{t_0}^{t} d\tau' \, f(\tau, \tau') \big[ A_S(\tau) A_S(\tau') C_0(\tau - \tau') \nonumber \\
&\quad + A_S'(\tau) A_S'(\tau') C_0^*(\tau - \tau') - 2 A_S(\tau) A_S'(\tau') C_0(\tau - \tau') \big]
\tag{13.11}
\end{align}

\textbf{Interpretation}: The factor $f(\tau, \tau')$ modulates the strength of correlations as the system evolves in cosmic time, implementing the $D = D_0 \cdot z$ gradient.
\end{derivation}

\section{Complete Derivation of the Correlation Function with Anisotropy and Gradient}

\begin{derivation}[Explicit Calculation of $\mathcal{C}(\tau, \tau')$ for the $T_\mu$ Field]
\textbf{Step 1: Definition of the Correlation Function}
\begin{equation}
\mathcal{C}(\tau, \tau') = \langle B(\tau) B(\tau') \rangle_T
\tag{15.1}
\end{equation}
with $B(\tau)$ given by (13.3).

\textbf{Step 2: Expansion of Reservoir Operators}
\begin{align}
B(\tau) &= \hbar \sum_{k,\lambda} \left[ g_{k}^{(\lambda)}(\tau) b_{k}^{(\lambda)} e^{-i\Omega_k(\tau)} + g_{k}^{(\lambda)*}(\tau) b_{k}^{(\lambda)\dagger} e^{i\Omega_k(\tau)} \right]
\tag{15.2} \\
B(\tau') &= \hbar \sum_{k',\lambda'} \left[ g_{k'}^{(\lambda')}(\tau') b_{k'}^{(\lambda')} e^{-i\Omega_{k'}(\tau')} + g_{k'}^{(\lambda')*}(\tau') b_{k'}^{(\lambda')\dagger} e^{i\Omega_{k'}(\tau')} \right]
\tag{15.3}
\end{align}

\textbf{Step 3: Operator Product}
\begin{align}
B(\tau)B(\tau') &= \hbar^2 \sum_{k,k'} \sum_{\lambda,\lambda'} \big[ g_{k}^{(\lambda)}(\tau) g_{k'}^{(\lambda')}(\tau') b_{k}^{(\lambda)} b_{k'}^{(\lambda')} e^{-i[\Omega_k(\tau)+\Omega_{k'}(\tau')]} \\
&\quad + g_{k}^{(\lambda)}(\tau) g_{k'}^{(\lambda')*}(\tau') b_{k}^{(\lambda)} b_{k'}^{(\lambda')\dagger} e^{-i[\Omega_k(\tau)-\Omega_{k'}(\tau')]} \\
&\quad + g_{k}^{(\lambda)*}(\tau) g_{k'}^{(\lambda')}(\tau') b_{k}^{(\lambda)\dagger} b_{k'}^{(\lambda')} e^{i[\Omega_k(\tau)-\Omega_{k'}(\tau')]} \\
&\quad + g_{k}^{(\lambda)*}(\tau) g_{k'}^{(\lambda')*}(\tau') b_{k}^{(\lambda)\dagger} b_{k'}^{(\lambda')\dagger} e^{i[\Omega_k(\tau)+\Omega_{k'}(\tau')]} \big]
\tag{15.4}
\end{align}

\textbf{Step 4: Thermal Expectation Values}
For the $T_\mu$ vector field in a thermal state:
\begin{align}
\langle b_{k}^{(\lambda)} b_{k'}^{(\lambda')} \rangle_T &= 0
\tag{15.5} \\
\langle b_{k}^{(\lambda)\dagger} b_{k'}^{(\lambda')\dagger} \rangle_T &= 0
\tag{15.6} \\
\langle b_{k}^{(\lambda)} b_{k'}^{(\lambda')\dagger} \rangle_T &= (n_k^{(\lambda)} + 1) \delta_{kk'} \delta_{\lambda\lambda'}
\tag{15.7} \\
\langle b_{k}^{(\lambda)\dagger} b_{k'}^{(\lambda')} \rangle_T &= n_k^{(\lambda)} \delta_{kk'} \delta_{\lambda\lambda'}
\tag{15.8}
\end{align}
where $n_k^{(\lambda)} = \frac{1}{e^{\beta\hbar\omega_k^{(\lambda)}} - 1}$ is the Bose-Einstein distribution for mode $(k,\lambda)$.

\textbf{Step 5: Non-Zero Terms}
Only terms with $b_{k}^{(\lambda)} b_{k'}^{(\lambda')\dagger}$ and $b_{k}^{(\lambda)\dagger} b_{k'}^{(\lambda')}$ contribute:
\begin{align}
\mathcal{C}(\tau, \tau') &= \hbar^2 \sum_{k,\lambda} \big[ g_{k}^{(\lambda)}(\tau) g_{k}^{(\lambda)*}(\tau') (n_k^{(\lambda)} + 1) e^{-i[\Omega_k(\tau)-\Omega_k(\tau')]} \\
&\quad + g_{k}^{(\lambda)*}(\tau) g_{k}^{(\lambda)}(\tau') n_k^{(\lambda)} e^{i[\Omega_k(\tau)-\Omega_k(\tau')]} \big]
\tag{15.9}
\end{align}

\textbf{Step 6: Simplified Form}
\begin{equation}
\mathcal{C}(\tau, \tau') = \hbar^2 \sum_{k,\lambda} |g_{k0}^{(\lambda)}|^2 f_k(\tau, \tau') \left[ (n_k^{(\lambda)} + 1) e^{-i[\Omega_k(\tau)-\Omega_k(\tau')]} + n_k^{(\lambda)} e^{i[\Omega_k(\tau)-\Omega_k(\tau')]} \right]
\tag{15.10}
\end{equation}
where $f_k(\tau, \tau') = (1+z(\tau))^{\gamma_k} (1+z(\tau'))^{\gamma_k}$ captures the gradient, and $g_{k0}^{(\lambda)}$ are the coefficients at $t=0$.

\textbf{Step 7: Low Frequency Limit (Cosmological Scales)}
For $\omega_k \ll H_0$, we can approximate $\Omega_k(\tau) \approx \omega_k \tau$ and $n_k^{(\lambda)} \approx \frac{k_B T_{\text{causal}}}{\hbar\omega_k} \gg 1$ (classical limit):
\begin{equation}
\mathcal{C}(\tau, \tau') \approx 2\hbar k_B T_{\text{causal}} \sum_{k,\lambda} \frac{|g_{k0}^{(\lambda)}|^2}{\omega_k} f_k(\tau, \tau') \cos[\omega_k (\tau - \tau')]
\tag{15.11}
\end{equation}
\end{derivation}

\section{Derivation of the Symmetric Component with Gradient}

\begin{derivation}[Noise and Decoherence Term with Temporal Evolution]
\textbf{Step 1: Definition of the Symmetric Component}
\begin{equation}
\mathcal{C}_{\text{sym}}(\tau, \tau') = \frac{1}{2} [\mathcal{C}(\tau, \tau') + \mathcal{C}^*(\tau, \tau')]
\tag{16.1}
\end{equation}

\textbf{Step 2: Complex Conjugate}
From (15.10):
\begin{equation}
\mathcal{C}^*(\tau, \tau') = \hbar^2 \sum_{k,\lambda} |g_{k0}^{(\lambda)}|^2 f_k(\tau, \tau') \left[ (n_k^{(\lambda)} + 1) e^{i[\Omega_k(\tau)-\Omega_k(\tau')]} + n_k^{(\lambda)} e^{-i[\Omega_k(\tau)-\Omega_k(\tau')]} \right]
\tag{16.2}
\end{equation}

\textbf{Step 3: Sum of Functions}
\begin{align}
\mathcal{C} + \mathcal{C}^* &= \hbar^2 \sum_{k,\lambda} |g_{k0}^{(\lambda)}|^2 f_k(\tau, \tau') \left[ (2n_k^{(\lambda)} + 1) e^{-i[\Omega_k(\tau)-\Omega_k(\tau')]} + (2n_k^{(\lambda)} + 1) e^{i[\Omega_k(\tau)-\Omega_k(\tau')]} \right] \\
&= 2\hbar^2 \sum_{k,\lambda} |g_{k0}^{(\lambda)}|^2 f_k(\tau, \tau') (2n_k^{(\lambda)} + 1) \cos[\Omega_k(\tau) - \Omega_k(\tau')]
\tag{16.3}
\end{align}

\textbf{Step 4: Approximation for Small Temporal Differences}
For $\Delta\tau = \tau - \tau'$ small compared to the Hubble time $H^{-1}$, we have $\Omega_k(\tau) - \Omega_k(\tau') \approx \omega_k(\bar{\tau}) \Delta\tau$, where $\bar{\tau} = (\tau+\tau')/2$. Thus:
\begin{equation}
\mathcal{C} + \mathcal{C}^* \approx 2\hbar^2 \sum_{k,\lambda} |g_{k0}^{(\lambda)}|^2 f_k(\tau, \tau') (2n_k^{(\lambda)} + 1) \cos[\omega_k(\bar{\tau}) \Delta\tau]
\tag{16.4}
\end{equation}

\textbf{Step 5: Final Symmetric Component}
\begin{equation}
\mathcal{C}_{\text{sym}}(\tau, \tau') = \hbar^2 \sum_{k,\lambda} |g_{k0}^{(\lambda)}|^2 f_k(\tau, \tau') (2n_k^{(\lambda)} + 1) \cos[\omega_k(\bar{\tau}) \Delta\tau]
\tag{16.5}
\end{equation}

\textbf{Step 6: Physical Interpretation}
$\mathcal{C}_{\text{sym}}$ represents the quantum and thermal noise causing decoerence. The factor $f_k(\tau, \tau') = (1+z(\tau))^{\gamma_k}(1+z(\tau'))^{\gamma_k}$ shows that noise strength \textbf{grows with redshift}, implementing the observed gradient. For $\gamma_k \approx 1$, we have $f_k(\tau, \tau') \propto (1+z)^2$, corresponding to $D \propto z$.
\end{derivation}

\section{Derivation of the Antisymmetric Component with Precession}

\begin{derivation}[Dissipation Term and Lamb Shift with Directional Dependence]
\textbf{Step 1: Definition of the Antisymmetric Component}
\begin{equation}
\mathcal{C}_{\text{skew}}(\tau, \tau') = \frac{1}{2i} [\mathcal{C}(\tau, \tau') - \mathcal{C}^*(\tau, \tau')]
\tag{17.1}
\end{equation}

\textbf{Step 2: Difference of Functions}
\begin{align}
\mathcal{C} - \mathcal{C}^* &= \hbar^2 \sum_{k,\lambda} |g_{k0}^{(\lambda)}|^2 f_k(\tau, \tau') \left[ (n_k^{(\lambda)}+1) e^{-i\Delta\Omega} + n_k^{(\lambda)} e^{i\Delta\Omega} - (n_k^{(\lambda)}+1) e^{i\Delta\Omega} - n_k^{(\lambda)} e^{-i\Delta\Omega} \right] \\
&= \hbar^2 \sum_{k,\lambda} |g_{k0}^{(\lambda)}|^2 f_k(\tau, \tau') \left[ e^{-i\Delta\Omega} - e^{i\Delta\Omega} \right] \\
&= -2i\hbar^2 \sum_{k,\lambda} |g_{k0}^{(\lambda)}|^2 f_k(\tau, \tau') \sin[\Delta\Omega]
\tag{17.2}
\end{align}
where $\Delta\Omega = \Omega_k(\tau) - \Omega_k(\tau')$.

\textbf{Step 3: Final Antisymmetric Component}
\begin{equation}
\mathcal{C}_{\text{skew}}(\tau, \tau') = -\hbar^2 \sum_{k,\lambda} |g_{k0}^{(\lambda)}|^2 f_k(\tau, \tau') \sin[\Omega_k(\tau) - \Omega_k(\tau')]
\tag{17.3}
\end{equation}

\textbf{Step 4: Inclusion of Precession}
The cosmic precession $\theta_T(z) = \theta_0 + \omega_p z$ modifies the coefficients through angular dependence. We can write:
\begin{equation}
|g_{k0}^{(\lambda)}|^2 \rightarrow |g_{k0}^{(\lambda)}|^2 \cdot [1 + \alpha \cos^2(\theta_k - \theta_T(z(\bar{\tau})))]
\tag{17.4}
\end{equation}
where $\alpha$ is the measured intrinsic anisotropy, $\theta_k$ is the direction of $\mathbf{k}$, and $\theta_T(z)$ evolves with the precessional parameter $\omega_p$.

\textbf{Step 5: Complete Form with Anisotropy and Precession}
\begin{equation}
\mathcal{C}_{\text{skew}}(\tau, \tau') = -\hbar^2 \sum_{k,\lambda} |g_{k0}^{(\lambda)}|^2 f_k(\tau, \tau') [1 + \alpha \cos^2(\theta_k - \theta_T(z(\bar{\tau})))] \sin[\omega_k(\bar{\tau}) \Delta\tau]
\tag{17.5}
\end{equation}

\textbf{Step 6: Physical Interpretation}
$\mathcal{C}_{\text{skew}}$ represents dissipative effects and contributes to the Lamb shift. The angular modulation $[1 + \alpha \cos^2(\theta_k - \theta_T)]$ introduces directional dependence in dissipation, while the precession $\theta_T(z(\bar{\tau}))$ causes this preferred direction to rotate with cosmic time.

\textbf{Connection to Anisotropic Linear Response}: The antisymmetric component is related to the reservoir's response function via generalized Kramers-Kronig relations for anisotropic media.
\end{derivation}

\section{Derivation of the Anisotropic Spectral Density with Evolution}

\begin{derivation}[Definition and Properties of $J(\omega, \theta, t)$ in RRT]
\textbf{Step 1: General Definition of Spectral Density}
\begin{equation}
J(\omega, \theta, t) = \sum_{k,\lambda} |g_{k}^{(\lambda)}(t)|^2 \delta(\omega - \omega_k(t)) \delta_{\theta, \theta_k}
\tag{18.1}
\end{equation}
where $\theta$ is the angular direction and we include explicit time dependence in $g_k$ and $\omega_k$.

\textbf{Step 2: Parametric Form with Gradient}
Using (8.6):
\begin{equation}
J(\omega, \theta, t) = (1+z(t))^{2\gamma} J_0(\omega, \theta)
\tag{18.2}
\end{equation}
where $J_0(\omega, \theta)$ is the spectral density at $t=0$.

\textbf{Step 3: Anisotropic Form}
From the analysis of the SDSS SuperSet v3, we know the spectral density is modulated by:
\begin{equation}
J_0(\omega, \theta) = J_{\text{iso}}(\omega) \left[ 1 + \alpha \cos^2(\theta - \theta_0) \right]
\end{equation}
where $\alpha$ represents the amplitude of the intrinsic dipole. The stability of the causal vector direction ($\theta_0$) was subjected to a 50-iteration Jackknife stress test, resulting in strict angular dispersion ($\pm 1.69^\circ$). This stabilized angular dispersion, associated with the \textbf{30.36$\sigma$}, proves that the precession phase is invariant to random data removal, shielding RRT against selection bias or local fluctuation hypotheses.

\textbf{Step 4: Inclusion of Precession}
The anisotropy direction evolves with redshift:
\begin{equation}
J(\omega, \theta, t) = (1+z(t))^{2\gamma} J_{\text{iso}}(\omega) \left[ 1 + \alpha \cos^2(\theta - \theta_T(z(t))) \right]
\tag{18.4}
\end{equation}
with $\theta_T(z) = \theta_0 + \omega_p z$, and $\omega_p$ is the Cosmic Precession Rate governed by the cosmic expansion.

\textbf{Step 5: Continuous Representation}
In the continuous limit:
\begin{equation}
J(\omega, \theta, t) = (1+z(t))^{2\gamma} \int d\omega' \, \rho(\omega') |g(\omega')|^2 \delta(\omega - \omega') \left[ 1 + \alpha \cos^2(\theta - \theta_T(z(t))) \right]
\tag{18.5}
\end{equation}
where $\rho(\omega)$ is the density of states.

\textbf{Step 6: Correlation Function in terms of $J(\omega, \theta, t)$}
\begin{align}
\mathcal{C}(\tau, \tau') &= \int_0^\infty d\omega \int d\Omega \, J(\omega, \theta, \bar{\tau}) \left[ (n(\omega, T_{\text{causal}}) + 1) e^{-i\omega(\tau-\tau')} + n(\omega, T_{\text{causal}}) e^{i\omega(\tau-\tau')} \right] \\
&\quad \times (1+z(\tau))^\gamma (1+z(\tau'))^\gamma \left[ 1 + \alpha \cos^2(\theta - \theta_T(z(\bar{\tau}))) \right]
\tag{18.6}
\end{align}
where $d\Omega$ is the solid angle element, $\bar{\tau} = (\tau+\tau')/2$, and $n(\omega, T) = (e^{\beta\hbar\omega}-1)^{-1}$.

\textbf{Step 7: Typical Forms of Isotropic Spectral Density}
\begin{itemize}
\item \textbf{Ohmic}: $J_{\text{iso}}(\omega) = \eta\omega e^{-\omega/\omega_c}$
\item \textbf{Super-Ohmic}: $J_{\text{iso}}(\omega) = \eta\omega^s e^{-\omega/\omega_c}$ with $s > 1$
\item \textbf{Sub-Ohmic}: $J_{\text{iso}}(\omega) = \eta\omega^s e^{-\omega/\omega_c}$ with $s < 1$
\end{itemize}
In RRT, cosmological data suggest $s \approx 1$ (Ohmic) for $\omega \ll \omega_c$, with $\omega_c \sim H_0$.

\textbf{Step 8: Cosmic Cut-off Frequency}
$\omega_c$ represents the maximum frequency scale of the causal reservoir, potentially related to the inverse Planck time or the inverse Hubble radius.

\textbf{Step 9: Final Spectral Density for Calculations}
For most cosmological applications, we apply the generalized form:
\begin{equation}
    J(\omega, \theta, t) = \eta \omega e^{-\omega/\omega_c} (1+z(t))^{2} \left[ 1 + \alpha \cos^2(\theta - \theta_T(z(t))) \right]
\tag{18.7}
\end{equation}
where the $(1+z)^2$ factor derives from the structural gradient $D = D_0 \cdot z$, and the Cortez Precession is encapsulated in the non-stationary angular argument.

\textbf{Connection to Cosmic Anisotropy}: The form (18.7) mathematically encapsulates all anisotropy observations: the linear gradient in $z$, the intrinsic anisotropy $\alpha$, the initial direction $\theta_0$, and the precession $\omega_p$. This is the fundamental bridge between open systems formalism and observational cosmology.

\textbf{Step 10: The Planck Verdict (Primordial Asymmetry)}
The analysis of low multipoles ($l < 20$) of the CMB reveals a \textbf{98.36\%} alignment with the Cortez Axis. In Phase 0 (Primordial Ocean), the absence of baryonic islands allowed $\beta$ viscosity to establish a pure preferred direction. The "Axis of Evil" is the scar left by the initial viscous flow, proving RRT as the universe's origin geometry.
\end{derivation}

This destructive interference explains why isotropic audits fail. The Cortez Rotation requires that the analysis be stratified. Without the phase correction, the statistical signal of \textbf{30.36$\sigma$} identified in the SDSS resonance stratum would be diluted into noise, as demonstrated by the marginal significance in noisy photometric catalogs.

% --- BLOCK 3: Canonical Expansion and Master Equation in RRT ---

\section*{BLOCK 3: Canonical Expansion and Master Equation in RRT}

\textbf{Conceptual Transition}: With the influence functional derived, we now move to the derivation of the master equation for the density matrix of the system coupled to the causal reservoir. The crucial novelty is that the dissipation and decoherence coefficients now possess temporal (due to the gradient) and angular (due to anisotropy and precession) dependence.

\section{Derivation of the Mean and Difference Basis with Cosmic Dependence}

\begin{derivation}[Variable Change in Path Space with Cosmic Time]
\textbf{Step 1: Variable Definition}
We introduce variables that separate the mean classical trajectory from quantum fluctuations, now with cosmic time awareness:
\begin{align}
A_m(\tau) &= \frac{1}{2} (A_S(\tau) + A_S'(\tau))
\tag{19.1} \\
\Delta A(\tau) &= A_S(\tau) - A_S'(\tau)
\tag{19.2}
\end{align}

\textbf{Step 2: Inverse Transformation}
\begin{align}
A_S(\tau) &= A_m(\tau) + \frac{1}{2} \Delta A(\tau)
\tag{19.3} \\
A_S'(\tau) &= A_m(\tau) - \frac{1}{2} \Delta A(\tau)
\tag{19.4}
\end{align}

\textbf{Step 3: Transformation Jacobian}
The transformation is linear with determinant 1; therefore, the Jacobian is unitary:
\begin{equation}
\mathscr{D}A_S \mathscr{D}A_S' = \mathscr{D}A_m \mathscr{D}\Delta A
\tag{19.5}
\end{equation}

\textbf{Step 4: Physical Interpretation in RRT}
\begin{itemize}
\item $A_m(\tau)$: Represents the mean classical trajectory of the system observable.
\item $\Delta A(\tau)$: Represents the quantum fluctuations (coherences) whose suppression defines decoherence.
\end{itemize}
In RRT, both may have implicit dependence on cosmic time through their evolution in the expanding universe.

\textbf{Step 5: Boundary Conditions for the Density Matrix}
For the evolution of the density matrix $\rho_S$, we have:
\begin{equation}
\Delta A(t_0) = \Delta A(t) = 0
\tag{19.6}
\end{equation}
while $A_m$ has free boundary conditions. These conditions reflect that in the path integral formalism for $\rho_S$, the "forward" and "backward" trajectories coincide at the initial and final times.
\end{derivation}

\section{Derivation of the Real Component \texorpdfstring{$\mathcal{I}_{\text{Re}}$}{I\_Re} with Gradient}

\begin{derivation}[Decoherence Term in the Canonical Basis with Temporal Evolution]
\textbf{Step 1: Original Functional Expression}
From (14.6) and (16.5), the real part of the exponent is:
\begin{equation}
\mathcal{I}_{\text{Re}} = -\frac{1}{2\hbar^2} \int_{t_0}^{t} d\tau \int_{t_0}^{t} d\tau' \, f(\tau,\tau') \mathcal{C}_{\text{sym}}(\tau,\tau') \left[ A_S(\tau)A_S(\tau') + A_S'(\tau)A_S'(\tau') - 2A_S(\tau)A_S'(\tau') \right]
\tag{20.1}
\end{equation}
with $f(\tau,\tau') = (1+z(\tau))^\gamma (1+z(\tau'))^\gamma$ and $\mathcal{C}_{\text{sym}}$ given by (16.5).

\textbf{Step 2: Substitution into the Canonical Basis}
Substituting $A_S = A_m + \frac{1}{2}\Delta A$ and $A_S' = A_m - \frac{1}{2}\Delta A$:

\textbf{Step 3: Term-by-Term Development}
We calculate each term separately:
\begin{align}
A_SA_S &= \left(A_m + \frac{1}{2}\Delta A\right)\left(A_m' + \frac{1}{2}\Delta A'\right) \\
&= A_mA_m' + \frac{1}{2}A_m\Delta A' + \frac{1}{2}\Delta A A_m' + \frac{1}{4}\Delta A\Delta A'
\tag{20.2}
\end{align}

\begin{align}
A_S'A_S' &= \left(A_m - \frac{1}{2}\Delta A\right)\left(A_m' - \frac{1}{2}\Delta A'\right) \\
&= A_mA_m' - \frac{1}{2}A_m\Delta A' - \frac{1}{2}\Delta A A_m' + \frac{1}{4}\Delta A\Delta A'
\tag{20.3}
\end{align}

\begin{align}
-2A_SA_S' &= -2\left(A_m + \frac{1}{2}\Delta A\right)\left(A_m' - \frac{1}{2}\Delta A'\right) \\
&= -2A_mA_m' + A_m\Delta A' - \Delta A A_m' + \frac{1}{2}\Delta A\Delta A'
\tag{20.4}
\end{align}

\textbf{Step 4: Sum of All Terms}
Summing (20.2), (20.3), and (20.4):
\begin{align}
[A_SA_S &+ A_S'A_S' - 2A_SA_S'] = \\
&\left[A_mA_m' + \frac{1}{2}A_m\Delta A' + \frac{1}{2}\Delta A A_m' + \frac{1}{4}\Delta A\Delta A'\right] + \\
&\left[A_mA_m' - \frac{1}{2}A_m\Delta A' - \frac{1}{2}\Delta A A_m' + \frac{1}{4}\Delta A\Delta A'\right] + \\
&\left[-2A_mA_m' + A_m\Delta A' - \Delta A A_m' + \frac{1}{2}\Delta A\Delta A'\right] \\
&= \Delta A(\tau)\Delta A(\tau')
\tag{20.5}
\end{align}
\textbf{Note}: The $A_m$ terms cancel out completely! This is a fundamental result: only quantum fluctuations $\Delta A$ contribute to decoherence.

\textbf{Step 5: Final Result for $\mathcal{I}_{\text{Re}}$}
\begin{equation}
\mathcal{I}_{\text{Re}} = -\frac{1}{2\hbar^2} \int_{t_0}^{t} d\tau \int_{t_0}^{t} d\tau' \, f(\tau,\tau') \mathcal{C}_{\text{sym}}(\tau,\tau') \Delta A(\tau) \Delta A(\tau')
\tag{20.6}
\end{equation}

\textbf{Step 6: Form with Spectral Density}
Substituting $\mathcal{C}_{\text{sym}}$ from (16.5):
\begin{equation}
\mathcal{I}_{\text{Re}} = -\frac{1}{2} \int_{t_0}^{t} d\tau \int_{t_0}^{t} d\tau' \, \Delta A(\tau)\Delta A(\tau') \int_0^\infty d\omega \, J(\omega,\theta,\bar{\tau}) (2n(\omega)+1) \cos[\omega(\tau-\tau')]
\tag{20.7}
\end{equation}
where $J(\omega,\theta,t)$ is given by (18.7) and $\bar{\tau} = (\tau+\tau')/2$.

\textbf{Interpretation}: Decoherence is caused by the accumulation of correlations between quantum fluctuations $\Delta A$ at different times, weighted by the reservoir's correlation function. The factor $f(\tau,\tau')$ causes the decoherence strength to \textbf{increase with redshift}, implementing the $D \propto z$ gradient.

\textbf{Updated Note}: With the new Golden Triad data, the $(1+z)^2$ evolution is confirmed with 0.5\% precision up to $z=2.5$.
\end{derivation}

\section{Derivation of the Imaginary Component \texorpdfstring{$\mathcal{I}_{\text{Im}}$}{I\_Im} with Precession}

\begin{derivation}[Dissipation Term in the Canonical Basis with Angular Dependence]
\textbf{Step 1: Original Imaginary Part Expression}
From (14.6) and (17.5):
\begin{equation}
i\mathcal{I}_{\text{Im}} = -\frac{1}{2\hbar^2} \int d\tau d\tau' \, f(\tau,\tau') i\mathcal{C}_{\text{skew}}(\tau,\tau') \left[ A_S(\tau)A_S(\tau') - A_S'(\tau)A_S'(\tau') - 2A_S(\tau)A_S'(\tau') \right]
\tag{21.1}
\end{equation}

\textbf{Step 2: Brackets Denotation}
We define $[\cdots]' = A_S(\tau)A_S(\tau') - A_S'(\tau)A_S'(\tau') - 2A_S(\tau)A_S'(\tau')$

\textbf{Step 3: Individual Terms Development}
First term:
\begin{align}
A_SA_S - A_S'A_S' &= \left(A_m + \frac{1}{2}\Delta A\right)\left(A_m' + \frac{1}{2}\Delta A'\right) - \left(A_m - \frac{1}{2}\Delta A\right)\left(A_m' - \frac{1}{2}\Delta A'\right) \\
&= A_m\Delta A' + \Delta A A_m'
\tag{21.2}
\end{align}

Second term:
\begin{align}
-2A_SA_S' &= -2\left(A_m + \frac{1}{2}\Delta A\right)\left(A_m' - \frac{1}{2}\Delta A'\right) \\
&= -2A_mA_m' + A_m\Delta A' - \Delta A A_m' + \frac{1}{2}\Delta A\Delta A'
\tag{21.3}
\end{align}

\textbf{Step 4: Full Sum}
\begin{align}
[\cdots]' &= [A_m\Delta A' + \Delta A A_m'] + [-2A_mA_m' + A_m\Delta A' - \Delta A A_m' + \frac{1}{2}\Delta A\Delta A'] \\
&= -2A_m(\tau)A_m(\tau') + 2A_m(\tau)\Delta A(\tau') + \frac{1}{2}\Delta A(\tau)\Delta A(\tau')
\tag{21.4}
\end{align}

\textbf{Step 5: Lindblad Approximation (Second Order)}
In the standard Lindblad master equation approximation:
\begin{itemize}
\item The $A_mA_m$ term is symmetric in $\tau \leftrightarrow \tau'$ and couples to $\mathcal{C}_{\text{skew}}$, which is antisymmetric $\rightarrow$ zero integral.
\item The $\Delta A\Delta A$ term is fourth-order in fluctuations $\rightarrow$ neglected in the second-order expansion.
\end{itemize}

Therefore:
\begin{equation}
[\cdots]' \approx 2A_m(\tau)\Delta A(\tau')
\tag{21.5}
\end{equation}

\footnote{Note: At very high redshifts ($z \gtrsim 10^3$), particularly in regimes such as the inflationary epoch, the condition $\tau_c \ll \tau_S$ may fail severely, invalidating the Lindblad approximation. In such cases, the exact Nakajima-Zwanzig master equation (Derivation 24) must be used directly, as it maintains the full integral memory of the non-Markovian kernel.}

\textbf{Step 6: Final Result for $\mathcal{I}_{\text{Im}}$}
\begin{equation}
i\mathcal{I}_{\text{Im}} \approx -\frac{i}{\hbar^2} \int_{t_0}^{t} d\tau \int_{t_0}^{t} d\tau' \, f(\tau,\tau') \mathcal{C}_{\text{skew}}(\tau,\tau') A_m(\tau) \Delta A(\tau')
\tag{21.6}
\end{equation}

\textbf{Step 7: Form with Anisotropic Spectral Density}
Substituting $\mathcal{C}_{\text{skew}}$ from (17.5):
\begin{align}
\mathcal{I}_{\text{Im}} &= -\frac{1}{\hbar^2} \int_{t_0}^{t} d\tau \int_{t_0}^{t} d\tau' \, A_m(\tau)\Delta A(\tau') \times \\
&\quad \int_0^\infty d\omega \, J(\omega,\theta,\bar{\tau}) [1 + \alpha \cos^2(\theta - \theta_T(z(\bar{\tau})))] \sin[\omega(\tau-\tau')]
\tag{21.7}
\end{align}

\textbf{Interpretation}: Dissipation couples the mean classical trajectory $A_m$ with quantum fluctuations $\Delta A$. The angular modulation $[1 + \alpha \cos^2(\theta - \theta_T)]$ introduces directional dependence, and the precession $\theta_T(z)$ causes this direction to rotate with cosmic time.

\textbf{Updated Note}: The value $\alpha = 1.198 \pm 0.015$ is now calibrated by the Golden Triad, representing a significant anisotropy in the causal coupling.
\end{derivation}

\section{Derivation of the Canonical Influence Functional in RRT}

\begin{derivation}[Final Form in the Mean and Difference Basis with Combined Effects]
\textbf{Step 1: Sum of Real and Imaginary Contributions}
\begin{align}
\ln \mathcal{F}[A_S, A_S'] &= \mathcal{I}_{\text{Re}} + i\mathcal{I}_{\text{Im}} \\
&= -\frac{1}{\hbar^2} \int_{t_0}^{t} d\tau \int_{t_0}^{t} d\tau' \, f(\tau,\tau') \big\{ \frac{1}{2} \mathcal{C}_{\text{sym}}(\tau,\tau') \Delta A(\tau) \Delta A(\tau') \\
&\quad + i \mathcal{C}_{\text{skew}}(\tau,\tau') A_m(\tau) \Delta A(\tau') \big\}
\tag{22.1}
\end{align}

\textbf{Step 2: Symmetry of the Double Integral}
We can symmetrize the imaginary term using the fact that $\mathcal{C}_{\text{skew}}(\tau,\tau')$ is antisymmetric:
\begin{align}
\int d\tau d\tau' \, \mathcal{C}_{\text{skew}}(\tau,\tau') A_m(\tau) \Delta A(\tau') &= \frac{1}{2} \int d\tau d\tau' \, \mathcal{C}_{\text{skew}}(\tau,\tau') [A_m(\tau) \Delta A(\tau') - A_m(\tau') \Delta A(\tau)]
\tag{22.2}
\end{align}

\textbf{Step 3: Final Symmetrized Form}
\begin{align}
\ln \mathcal{F}[A_S, A_S'] &= -\frac{1}{\hbar^2} \int_{t_0}^{t} d\tau \int_{t_0}^{t} d\tau' \, f(\tau,\tau') \big\{ \frac{1}{2} \mathcal{C}_{\text{sym}}(\tau,\tau') \Delta A(\tau) \Delta A(\tau') \\
&\quad + \frac{i}{2} \mathcal{C}_{\text{skew}}(\tau,\tau') [A_m(\tau) \Delta A(\tau') - A_m(\tau') \Delta A(\tau)] \big\}
\tag{22.3}
\end{align}

\textbf{Step 4: Full Physical Interpretation}
\begin{itemize}
\item \textbf{First term}: Pure decoherence - depends only on quantum fluctuations $\Delta A$. It suppresses off-diagonal coherences in the density matrix. The factor $f(\tau,\tau')$ causes the decoherence rate to increase with redshift.

\item \textbf{Second term}: Dissipation - couples the classical trajectory $A_m$ with fluctuations $\Delta A$. It causes relaxation toward equilibrium and the Lamb shift (energy renormalization). The angular dependence and precession are contained within $\mathcal{C}_{\text{skew}}$.
\end{itemize}

\textbf{Step 5: Markovian Limit with Gradient}
In the traditional Markovian limit, $\mathcal{C}_{\text{sym}}(\tau,\tau') \rightarrow \gamma \delta(\tau-\tau')$ and $\mathcal{C}_{\text{skew}}(\tau,\tau') \rightarrow \lambda \delta'(\tau-\tau')$. In RRT with a gradient, these become:
\begin{align}
\mathcal{C}_{\text{sym}}(\tau,\tau') &\rightarrow \gamma(z(\tau)) \delta(\tau-\tau')
\tag{22.4} \\
\mathcal{C}_{\text{skew}}(\tau,\tau') &\rightarrow \lambda(z(\tau)) \delta'(\tau-\tau')
\tag{22.5}
\end{align}
where $\gamma(z) = \gamma_0 (1+z)^{2\gamma}$ and $\lambda(z) = \lambda_0 (1+z)^{2\gamma}$ grow with redshift.

\textbf{Robustness Note (Jackknife Test)}: The calibration of coefficients via Jackknife resampling in SDSS DR16Q (after $3\sigma$ clipping, $N=216,076$) confirmed the signal stability with a directional standard deviation of only $1.69^\circ$. The consolidated value $D_0 = -0.0189 \pm 0.0005$ establishes the intensity of the causal gradient with statistical significance that shields the theory against local fluctuations, galactic dust, or selection bias.
\end{derivation}

\section{Derivation of the Keldysh Master Equation with Cosmic Time}

\begin{derivation}[Path Integral Formulation for the Density Matrix in RRT]
\textbf{Step 1: Path Integral Representation of the Density Matrix}
The reduced density matrix of the system evolves as:
\begin{equation}
\rho_S(t) = \int \mathscr{D}A_S \mathscr{D}A_S' \, e^{\frac{i}{\hbar}(S_S[A_S] - S_S[A_S'])} \mathcal{F}[A_S, A_S'] \rho_S(t_0)
\tag{23.1}
\end{equation}
where $S_S$ is the action of the isolated system and $\mathcal{F}$ is the derived influence functional.

\textbf{Step 2: Change to $A_m$ and $\Delta A$ Variables}
\begin{equation}
\rho_S(t) = \int \mathscr{D}A_m \mathscr{D}\Delta A \, e^{\frac{i}{\hbar}\int d\tau \, \Delta A(\tau) \frac{\delta S_S}{\delta A_m(\tau)} + \cdots} \mathcal{F}[A_m, \Delta A] \rho_S(t_0)
\tag{23.2}
\end{equation}
where we expand $S_S[A_S] - S_S[A_S']$ to first order in $\Delta A$.

\textbf{Step 3: Temporal Variation of the Density Matrix}
\begin{equation}
\frac{d\rho_S}{dt} = \lim_{\Delta t \to 0} \frac{\rho_S(t + \Delta t) - \rho_S(t)}{\Delta t}
\tag{23.3}
\end{equation}

\textbf{Step 4: Expansion for Small $\Delta t$}
For small $\Delta t$ (but still much larger than the reservoir correlation time $\tau_c$), we expand the influence functional:
\begin{equation}
\mathcal{F}[A_S, A_S'] \approx 1 - \frac{1}{\hbar^2} \int_t^{t+\Delta t} d\tau \int_t^{t+\Delta t} d\tau' \, f(\tau,\tau') [\cdots]
\tag{23.4}
\end{equation}

\textbf{Step 5: Dominant Term}
The dominant term arises from the double integral in the interval $[t, t+\Delta t]$:
\begin{equation}
\rho_S(t+\Delta t) \approx \rho_S(t) - \frac{1}{\hbar^2} \int_t^{t+\Delta t} d\tau \int_t^{t+\Delta t} d\tau' \, f(\tau,\tau') [\cdots] \rho_S(t)
\tag{23.5}
\end{equation}

\textbf{Step 6: Differential Equation}
\begin{equation}
\frac{d\rho_S}{dt} = -\frac{1}{\hbar^2} \lim_{\Delta t \to 0} \frac{1}{\Delta t} \int_t^{t+\Delta t} d\tau \int_t^{t+\Delta t} d\tau' \, f(\tau,\tau') [\cdots] \rho_S(t)
\tag{23.6}
\end{equation}
where $[\cdots]$ represents the appropriate combination of operators according to (22.3).

\textbf{Step 7: Operational Form}
In terms of system operators:
\begin{equation}
\frac{d\rho_S}{dt} = -\frac{i}{\hbar}[H_S, \rho_S] + \mathcal{L}_{\text{diss}}(z(t))[\rho_S]
\tag{23.7}
\end{equation}
where $\mathcal{L}_{\text{diss}}(z)$ is the dissipative superoperator that depends on redshift (via $f(\tau,\tau')$) and contains both decoherence and dissipation.

\textbf{Updated Note}: The $z(t)$ dependence is now empirically calibrated by the Golden Triad, yielding $\mathcal{L}_{\text{diss}}(z) = \mathcal{L}_{\text{diss}}^{(0)} (1+z)^{2.00 \pm 0.02}$.
\end{derivation}

\section{Derivation of the Markov Approximation with Gradient}

\begin{derivation}[Short Memory Limit with Slow Evolution]
\textbf{Step 1: Markovian Hypothesis in RRT}
We assume that reservoir correlations decay rapidly compared to the system's timescale, but allow for slow evolution of parameters due to cosmic expansion:
\begin{equation}
\mathcal{C}(\tau, \tau') \approx \Gamma(z(\bar{\tau})) \delta(\tau-\tau') + i\Lambda(z(\bar{\tau})) \delta'(\tau-\tau')
\tag{24.1}
\end{equation}
where $\bar{\tau} = (\tau+\tau')/2$, and $\Gamma(z)$, $\Lambda(z)$ evolve with $z$.

\textbf{Step 2: Moment Expansion with Evolution}
We can expand the correlation functions:
\begin{align}
\mathcal{C}_{\text{sym}}(\tau,\tau') &= \gamma_0(z(\bar{\tau})) \delta(\tau-\tau') + \gamma_1(z(\bar{\tau})) \delta'(\tau-\tau') + \cdots
\tag{24.2} \\
\mathcal{C}_{\text{skew}}(\tau,\tau') &= \lambda_0(z(\bar{\tau})) \delta(\tau-\tau') + \lambda_1(z(\bar{\tau})) \delta'(\tau-\tau') + \cdots
\tag{24.3}
\end{align}

\textbf{Step 3: Short Memory Approximation}
We keep only the dominant terms, but now with $z$-dependence:
\begin{align}
\mathcal{C}_{\text{sym}}(\tau,\tau') &\approx \gamma(z(\bar{\tau})) \delta(\tau-\tau')
\tag{24.4} \\
\mathcal{C}_{\text{skew}}(\tau,\tau') &\approx \lambda(z(\bar{\tau})) \delta'(\tau-\tau')
\tag{24.5}
\end{align}

\textbf{Step 4: Physical Justification}
This approximation is valid when:
\begin{equation}
\tau_c \ll \tau_S, \quad \text{and} \quad \left|\frac{1}{\gamma}\frac{d\gamma}{dt}\right| \tau_c \ll 1
\tag{24.6}
\end{equation}
where $\tau_c$ is the reservoir correlation time, $\tau_S$ is the system's characteristic time, and the second condition ensures that the evolution of $\gamma(z(t))$ is slow compared to $\tau_c$.

\textbf{Step 5: Effective Coefficients with Gradient}
\begin{align}
\gamma(z) &= \int_{-\infty}^{\infty} d\Delta\tau \, \mathcal{C}_{\text{sym}}(\bar{\tau}+\Delta\tau/2, \bar{\tau}-\Delta\tau/2) \\
&= \gamma_0 (1+z)^{2\gamma} \quad \text{with } \gamma_0 = \int_{-\infty}^{\infty} d\Delta\tau \, C_0(\Delta\tau)
\tag{24.7}
\end{align}

\begin{align}
\lambda(z) &= \int_{-\infty}^{\infty} d\Delta\tau \, \Delta\tau \, \mathcal{C}_{\text{skew}}(\bar{\tau}+\Delta\tau/2, \bar{\tau}-\Delta\tau/2) \\
&= \lambda_0 (1+z)^{2\gamma} \quad \text{with } \lambda_0 = \int_{-\infty}^{\infty} d\Delta\tau \, \Delta\tau \, C_0^{\text{skew}}(\Delta\tau)
\tag{24.8}
\end{align}

\textbf{Step 6: Validity for the Real Universe}
For the causal reservoir at cosmological scales:
\begin{itemize}
\item $\tau_c \sim \omega_c^{-1} \sim H_0^{-1} \sim 10^{17}$ s (Hubble time)
\item $\tau_S$ for typical quantum systems: $10^{-15}$ to $10^3$ s
\item $\left|\frac{1}{\gamma}\frac{d\gamma}{dt}\right| \sim H_0 \sim 10^{-18}$ s$^{-1}$
\end{itemize}
The condition $\tau_c \ll \tau_S$ \textbf{is not satisfied} for systems with $\tau_S < H_0^{-1}$! This means that the standard Markovian approximation may not be valid for coupling with the causal reservoir at cosmological scales. A non-Markovian treatment is required.

\textbf{Updated Note}: For astrophysical systems ($\tau_S \sim 10^{6}-10^{17}$ s) the Markovian approximation is marginally valid. For laboratory systems ($\tau_S < 10^3$ s) non-Markovian treatment is essential.
\end{derivation}

\section{Derivation of the Non-Markovian Dissipative Superoperator}

\begin{derivation}[Exact Form via Nakajima-Zwanzig Projection in RRT]
\textbf{Step 1: Nakajima-Zwanzig Projector}
We define the projector that traces over the reservoir:
\begin{equation}
\mathcal{P} \rho = \rho_T \otimes \Tr_T[\rho]
\tag{25.1}
\end{equation}
where $\rho_T$ is the state of the $T_\mu$ field (thermal with temperature $T_{\text{causal}}$), and $\mathcal{Q} = 1 - \mathcal{P}$.

\textbf{Step 2: Exact Projected Equation of Motion}
\begin{equation}
\frac{d}{dt} \mathcal{P}\rho(t) = -\frac{i}{\hbar} \mathcal{P}L \mathcal{P}\rho(t) - \frac{1}{\hbar^2} \int_0^t ds \, \mathcal{P}L e^{-i\mathcal{Q}L s/\hbar} \mathcal{Q}L \mathcal{P}\rho(t-s)
\tag{25.2}
\end{equation}
where $L = [H_{\text{total}}, \cdot]$ is the total Liouvillian, with $H_{\text{total}} = H_S + H_T + H_{\text{int}}$.

\textbf{Step 3: Born Approximation (Weak Coupling)}
We assume $\mathcal{P}\rho(t-s) \approx \mathcal{P}\rho(t)$ (the reservoir is not significantly perturbed):
\begin{equation}
\frac{d}{dt} \mathcal{P}\rho(t) \approx -\frac{i}{\hbar} \mathcal{P}L \mathcal{P}\rho(t) - \frac{1}{\hbar^2} \int_0^t ds \, \mathcal{P}L e^{-i\mathcal{Q}L_0 s/\hbar} \mathcal{Q}L \mathcal{P}\rho(t)
\tag{25.3}
\end{equation}
where $L_0 = [H_S + H_T, \cdot]$ is the free Liouvillian.

\textbf{Step 4: Expression for the Reduced Density Matrix}
For $\rho_S(t) = \Tr_T[\rho(t)]$, we obtain:
\begin{equation}
\frac{d\rho_S}{dt} = -\frac{i}{\hbar}[H_S, \rho_S] - \frac{1}{\hbar^2} \int_0^t d\tau \, \Tr_T \left[ H_{\text{int}}, [H_{\text{int}}(-\tau), \rho_S(t) \otimes \rho_T] \right]
\tag{25.4}
\end{equation}
where $H_{\text{int}}(-\tau) = e^{-i(H_S+H_T)\tau/\hbar} H_{\text{int}} e^{i(H_S+H_T)\tau/\hbar}$.

\textbf{Step 5: Commutator Development}
\begin{align}
\left[ H_{\text{int}}, [H_{\text{int}}(-\tau), \rho_S \otimes \rho_T] \right] &= H_{\text{int}} H_{\text{int}}(-\tau) \rho_S \rho_T - H_{\text{int}} \rho_S \rho_T H_{\text{int}}(-\tau) \\
&\quad - H_{\text{int}}(-\tau) \rho_S \rho_T H_{\text{int}} + \rho_S \rho_T H_{\text{int}}(-\tau) H_{\text{int}}
\tag{25.5}
\end{align}

\textbf{Step 6: Form with Explicit Time Dependence}
For $H_{\text{int}} = A_S \otimes B_T$ (where $B_T$ is the $T_\mu$ field operator):
\begin{align}
\frac{d\rho_S}{dt} &= -\frac{i}{\hbar}[H_S, \rho_S] - \frac{1}{\hbar^2} \int_0^t d\tau \, \big\{ C(\tau) [A_S, A_S(-\tau)\rho_S] \\
&\quad - C^*(-\tau) [A_S, \rho_S A_S(-\tau)] + \text{h.c.} \big\}
\tag{25.6}
\end{align}
where $C(\tau) = \langle B_T B_T(\tau) \rangle_T$ is the reservoir's correlation function, and $A_S(\tau) = e^{iH_S\tau/\hbar} A_S e^{-iH_S\tau/\hbar}$.

\textbf{Step 7: Inclusion of Gradient and Precession}
In RRT, $C(\tau)$ is non-stationary:
\begin{equation}
C(\tau; t) = f(t, t-\tau) C_0(\tau) \left[ 1 + \alpha \cos^2(\theta - \theta_T(z(t-\tau/2))) \right]
\tag{25.7}
\end{equation}
where $f(t, t-\tau) = (1+z(t))^\gamma (1+z(t-\tau))^\gamma$, and the dependence on $t$ enters explicitly.

\textbf{Updated Note}: The parameters $\alpha = 1.198$, and the precessional mechanism $\omega_p(z)$ are calibrated by the Golden Triad, ensuring that the non-Markovian equation correctly incorporates the observed cosmic evolution.
\end{derivation}

\section{Derivation of the Lindblad Form with Time-Dependent Coefficients}

\begin{derivation}[Representation of Jump Operators with Cosmic Evolution]
\textbf{Step 1: Master Equation in Non-Stationary Lindblad Form}
Even outside the strict Markovian regime, we can write a Lindblad form with time-dependent coefficients:
\begin{align}
\frac{d\rho_S}{dt} &= -\frac{i}{\hbar}[H_S + H_{LS}(t), \rho_S] \\
&\quad + \sum_{\omega} \sum_{\alpha\beta} \gamma_{\alpha\beta}(\omega, t) \left[ A_\beta(\omega) \rho_S A_\alpha^\dagger(\omega) - \frac{1}{2} \{ A_\alpha^\dagger(\omega) A_\beta(\omega), \rho_S \} \right]
\tag{25.1}
\end{align}

\textbf{Step 2: Spectral Operators}
The operators $A_\alpha(\omega)$ are defined in the energy basis of $H_S$:
\begin{equation}
A_\alpha(\omega) = \sum_{\epsilon' - \epsilon = \hbar\omega} \Pi(\epsilon) A_\alpha \Pi(\epsilon')
\tag{25.2}
\end{equation}
where $\Pi(\epsilon)$ is the projector onto the energy eigenstate $\epsilon$, and $A_\alpha$ are the coupling operators (in our case, typically $A_S$).

\textbf{Step 3: Time-Dependent Decay Coefficients}
\begin{equation}
\gamma_{\alpha\beta}(\omega, t) = \int_{-\infty}^{\infty} d\tau \, e^{i\omega\tau} C_{\alpha\beta}(\tau; t)
\tag{25.3}
\end{equation}
where $C_{\alpha\beta}(\tau; t)$ are the reservoir's correlation functions, now with explicit $t$-dependence due to the gradient.

\textbf{Step 4: Time-Dependent Lamb Shift Hamiltonian}
\begin{equation}
H_{LS}(t) = \sum_{\omega} \sum_{\alpha\beta} S_{\alpha\beta}(\omega, t) A_\alpha^\dagger(\omega) A_\beta(\omega)
\tag{25.4}
\end{equation}

with:
\begin{equation}
S_{\alpha\beta}(\omega, t) = \frac{1}{2i} \left[ \Gamma_{\alpha\beta}(\omega, t) - \Gamma_{\beta\alpha}^*(\omega, t) \right]
\tag{25.5}
\end{equation}
where $\Gamma_{\alpha\beta}(\omega, t) = \frac{1}{\hbar^2} \int_0^\infty d\tau \, e^{i\omega\tau} C_{\alpha\beta}(\tau; t)$.

\textbf{Step 5: Specific Form for RRT}
For $H_{\text{int}} = A_S \otimes B_T$ coupling:
\begin{align}
\gamma(\omega, t) &= \frac{1}{\hbar^2} \int_{-\infty}^{\infty} d\tau \, e^{i\omega\tau} f(t, t-\tau) C_0(\tau) \left[ 1 + \alpha \cos^2(\theta - \theta_T(z(t-\tau/2))) \right]
\tag{25.6} \\
&= \gamma_0(\omega) (1+z(t))^{2\gamma} \left[ 1 + \alpha \cos^2(\theta - \theta_T(z(t))) \right] \times F(\omega, t)
\tag{25.7}
\end{align}
where $F(\omega, t)$ is a form factor capturing deviations from perfect Markovianity.

\textbf{Step 6: Completeness Condition and Trace Preservation}
The Lindblad form automatically preserves:
\begin{align}
\Tr\left[\frac{d\rho_S}{dt}\right] &= 0
\tag{25.8} \\
\rho_S(t) &\succeq 0 \quad \forall t \quad \text{if} \quad \rho_S(0) \succeq 0
\tag{25.9}
\end{align}
even with time-dependent coefficients, provided that the matrix $\gamma_{\alpha\beta}(\omega, t)$ is positive semi-definite for all $t$.

\textbf{Step 7: Explicit Matrix Representation}
In the energy basis $\{|n\rangle\}$ of $H_S$:
\begin{align}
\frac{d\rho_{nm}}{dt} &= -i\omega_{nm}\rho_{nm} + \sum_{k,l} \mathcal{L}_{nm,kl}(t) \rho_{kl}
\tag{25.10}
\end{align}
where the time-dependent dissipative superoperator is:
\begin{equation}
\mathcal{L}_{nm,kl}(t) = \sum_{\omega} \gamma(\omega, t) \left[ A(\omega)_{nk} A^\dagger(\omega)_{lm} - \frac{1}{2} \delta_{ml} \sum_p A^\dagger(\omega)_{np} A(\omega)_{pk} - \frac{1}{2} \delta_{nk} \sum_p A^\dagger(\omega)_{lp} A(\omega)_{pm} \right]
\tag{25.11}
\end{equation}

\textbf{Step 8: Symmetries and Conservation Laws}
If the system possesses a conserved quantity $Q$ with $[H_S, Q] = 0$, then:
\begin{equation}
\frac{d\langle Q \rangle}{dt} = \Tr\left[Q \frac{d\rho_S}{dt}\right] = 0
\tag{25.12}
\end{equation}
provided that $[A_\alpha(\omega), Q] = 0$ for all $\alpha, \omega$, or weaker conditions are satisfied.
\end{derivation}

\section{Derivation of Dissipation Coefficients with Gradient}

\begin{derivation}[Explicit Calculation of $\gamma_{nm}(t)$ and $S_{nm}(t)$ in RRT]
\textbf{Step 1: Definition of the Time-Dependent Complex Integral}
We start from the general expression:
\begin{equation}
\Gamma_{nm}(t) = \frac{1}{\hbar^2} \int_0^\infty d\tau \, e^{i\omega_{nm}\tau} C(\tau; t)
\tag{26.1}
\end{equation}
with $C(\tau; t)$ given by (24.7).

\textbf{Step 2: Decomposition into Real and Imaginary Parts}
\begin{align}
\gamma_{nm}(t) &= 2\text{Re}[\Gamma_{nm}(t)] = \frac{2}{\hbar^2} \int_0^\infty d\tau \, \cos(\omega_{nm}\tau) C_{\text{sym}}(\tau; t)
\tag{26.2} \\
S_{nm}(t) &= \text{Im}[\Gamma_{nm}(t)] = \frac{1}{\hbar^2} \int_0^\infty d\tau \, \sin(\omega_{nm}\tau) C_{\text{skew}}(\tau; t)
\tag{26.3}
\end{align}

\textbf{Step 3: Substitution of RRT Correlation Functions}
Using (15.5) and (16.5):
\begin{align}
\gamma_{nm}(t) &= \frac{2}{\hbar^2} \int_0^\infty d\tau \, \cos(\omega_{nm}\tau) \sum_{k,\lambda} |g_{k0}^{(\lambda)}|^2 f_k(t, t-\tau) (2n_k^{(\lambda)}+1) \cos[\omega_k \tau]
\tag{26.4} \\
S_{nm}(t) &= -\frac{1}{\hbar^2} \int_0^\infty d\tau \, \sin(\omega_{nm}\tau) \sum_{k,\lambda} |g_{k0}^{(\lambda)}|^2 f_k(t, t-\tau) [1+\alpha\cos^2(\theta_k-\theta_T)] \sin[\omega_k \tau]
\tag{26.5}
\end{align}

\textbf{Step 4: Slow Evolution Approximation}
If the evolution of $f_k(t, t-\tau)$ is slow compared to $1/\omega_{nm}$, we can approximate $f_k(t, t-\tau) \approx f_k(t, t) = (1+z(t))^{2\gamma_k}$. Thus:
\begin{align}
\gamma_{nm}(t) &\approx (1+z(t))^{2\gamma} \cdot \frac{2}{\hbar^2} \int_0^\infty d\tau \, \cos(\omega_{nm}\tau) \sum_{k,\lambda} |g_{k0}^{(\lambda)}|^2 (2n_k^{(\lambda)}+1) \cos[\omega_k \tau]
\tag{26.6} \\
&= (1+z(t))^{2\gamma} \cdot \gamma_{nm}^{(0)}
\end{align}
where $\gamma_{nm}^{(0)}$ are the coefficients at $t=0$ ($z=0$).

\textbf{Step 5: Use of Integral Identities}
We apply the relations:
\begin{align}
\int_0^\infty d\tau \, \cos(\omega\tau)\cos(\omega'\tau) &= \frac{\pi}{2} [\delta(\omega-\omega') + \delta(\omega+\omega Pip)]
\tag{26.7} \\
\int_0^\infty d\tau \, \sin(\omega\tau)\sin(\omega'\tau) &= \frac{\pi}{2} [\delta(\omega-\omega') - \delta(\omega+\omega Pip)]
\tag{26.8}
\end{align}

\textbf{Step 6: Final Results for $\gamma_{nm}(t)$}
\begin{equation}
\gamma_{nm}(t) = (1+z(t))^{2\gamma} \cdot \frac{\pi}{\hbar^2} J(|\omega_{nm}|, \theta, t) [2n(|\omega_{nm}|, T_{\text{causal}}) + 1 + \text{sgn}(\omega_{nm})]
\tag{26.9}
\end{equation}
where $J(\omega, \theta, t)$ is given by (17.7), and $\text{sgn}(\omega_{nm}) = +1$ for $\omega_{nm} > 0$, and $-1$ for $\omega_{nm} < 0$.

\textbf{Step 7: Final Results for $S_{nm}(t)$ (Lamb shift)}
\begin{equation}
S_{nm}(t) = (1+z(t))^{2\gamma} \cdot \frac{1}{\hbar^2} \mathcal{P} \int_0^\infty d\omega \, J(\omega, \theta, t) \left[ \frac{2n(\omega, T_{\text{causal}})+1}{\omega_{nm}-\omega} - \frac{2n(\omega, T_{\text{causal}})+1}{\omega_{nm}+\omega} \right]
\tag{26.10}
\end{equation}
where $\mathcal{P}$ denotes the Cauchy principal value.

\textbf{Step 8: Important Limiting Cases}
\begin{itemize}
\item \textbf{Zero temperature} ($T_{\text{causal}} \to 0$, $n(\omega) \to 0$):
\begin{align}
\gamma_{nm}(t) &= (1+z(t))^{2\gamma} \cdot \frac{\pi}{\hbar^2} J(|\omega_{nm}|, \theta, t) [1 + \text{sgn}(\omega_{nm})]
\tag{26.11} \\
S_{nm}(t) &= (1+z(t))^{2\gamma} \cdot \frac{1}{\hbar^2} \mathcal{P} \int_0^\infty d\omega \, J(\omega, \theta, t) \left[ \frac{1}{\omega_{nm}-\omega} - \frac{1}{\omega_{nm}+\omega} \right]
\tag{26.12}
\end{align}

\item \textbf{High temperatures} ($k_B T_{\text{causal}} \gg \hbar\omega_{nm}$):
\begin{equation}
\gamma_{nm}(t) \approx (1+z(t))^{2\gamma} \cdot \frac{2\pi k_B T_{\text{causal}}}{\hbar^2 \omega_{nm}} J(|\omega_{nm}|, \theta, t)
\tag{26.13}
\end{equation}
\end{itemize}

\textbf{Step 9: Generalized Kramers-Kronig Relations}
The dissipative and reactive coefficients satisfy:
\begin{equation}
S_{nm}(t) = \frac{1}{\pi} \mathcal{P} \int_{-\infty}^\infty d\omega' \, \frac{\gamma_{nm}(\omega', t)}{\omega_{nm} - \omega'}
\tag{26.14}
\end{equation}
even with time dependence, provided the evolution is slow compared to the relevant frequency scales.

\textbf{Step 10: Physical Interpretation}
$\gamma_{nm}(t)$ represents transition rates between states $|n\rangle$ and $|m\rangle$ that \textbf{grow with redshift} as $(1+z)^{2\gamma}$. $S_{nm}(t)$ describes energy shifts (Lamb shift) that also increase with $z$. The angular dependence in $J(\omega,\theta,t)$ introduces anisotropy into the rates, and the precession $\theta_T(z(t))$ causes this anisotropy to rotate with cosmic time.

\textbf{Updated Note}: Calibration with the Golden Triad yields $\gamma = 1.00 \pm 0.02$, confirming the $(1+z)^2$ law with high precision.
\end{derivation}

\section{Derivation of the Final RRT Causal Master Equation}

\begin{derivation}[Explicit Form for Referential Relativity Theory]
\textbf{Step 1: System Operators in RRT}
For RRT, the interaction operator is typically $A_S$ (matter density, momentum, spin, etc.) with matrix elements $A_{nm} = \langle n|A_S|m\rangle$ in the energy basis of $H_S$.

\textbf{Step 2: Specific Transition Rates with Gradient}
From (26.9), with $\gamma=1$ to reproduce $D = D_0 \cdot z$:
\begin{equation}
\gamma_{nm}(t) = (1+z(t))^{2} \cdot \frac{\pi}{\hbar^2} J(|\omega_{nm}|, \theta, t) [2n(|\omega_{nm}|, T_{\text{causal}}) + 1 + \text{sgn}(\omega_{nm})]
\tag{27.1}
\end{equation}

\textbf{Step 3: Causal Lamb Shift with Evolution}
\begin{equation}
H_{LS}(t) = \sum_{n,m} \hbar S_{nm}(t) A_{nm} A_{mn} |n\rangle\langle n|
\tag{27.2}
\end{equation}
with $S_{nm}(t)$ given by (26.10).

\textbf{Step 4: Final Causal Lindblad Master Equation}
\begin{align}
\frac{d\rho_S}{dt} &= -i[H_S + H_{LS}(t), \rho_S] \\
&\quad + \sum_{n,m} \gamma_{nm}(t) \left( A_{nm} \rho_S A_{mn} - \frac{1}{2} \{ A_{mn} A_{nm}, \rho_S \} \right)
\tag{27.3}
\end{align}

\textbf{Step 5: Total Effective Hamiltonian}
\begin{equation}
H_{\text{eff}}(t) = H_S + H_{LS}(t)
\tag{27.4}
\end{equation}

\textbf{Step 6: Conservation Properties}
The equation preserves:
\begin{itemize}
\item \textbf{Trace}: $\Tr[\rho_S(t)] = 1$ for all $t$.
\item \textbf{Positivity}: $\rho_S(t) \succeq 0$ if $\rho_S(0) \succeq 0$.
\item \textbf{Mean Energy}: $\frac{d}{dt}\langle H_S \rangle = \Tr[\mathcal{L}(t)(\rho_S) H_S] = \sum_{n,m} \gamma_{nm}(t) \hbar\omega_{nm} |A_{nm}|^2 \rho_{nn}$.
\end{itemize}

\textbf{Step 7: Formal Solution}
The solution can be written as:
\begin{equation}
\rho_S(t) = \mathcal{T} \exp\left[ \int_{t_0}^t d\tau \, \mathcal{L}(\tau) \right] \rho_S(t_0)
\tag{27.5}
\end{equation}
where $\mathcal{T}$ is the time-ordering operator (required because $\mathcal{L}(t)$ does not commute with itself at different times), and $\mathcal{L}(t)$ is the time-dependent generator given by the right-hand side of (27.3).

\textbf{Step 8: Adiabatic Evolution Approximation}
If the evolution of $\gamma_{nm}(t)$ and $S_{nm}(t)$ is slow compared to the decoherence rates $1/\gamma_{nm}$, we can approximate:
\begin{equation}
\rho_S(t) \approx \exp\left[ \mathcal{L}(t) (t-t_0) \right] \rho_S(t_0)
\tag{27.6}
\end{equation}
where $\mathcal{L}(t)$ is evaluated at the current time $t$.

\textbf{Step 9: Connection to Cosmological Observations}
The parameters in the master equation are linked to observables:
\begin{itemize}
\item $(1+z(t))^2$ comes from the gradient $D = D_0 \cdot z$.
\item $J(\omega,\theta,t)$ contains $\alpha=1.198$ and direction $\theta_0=(125.5^\circ,-15.0^\circ)$.
\item $\theta_T(z) = \theta_0 + \omega_p z$ (Cortez Rotation).
\item $T_{\text{causal}}$ is the effective temperature of the $T_\mu$ field.
\end{itemize}

\textbf{Step 10: Limits of Applicability}
Equation (27.3) is valid for:
\begin{itemize}
\item Systems with $\tau_S \gg \tau_{\text{micro}}$ (microscopic scale of the system).
\item Weak coupling: $|\gamma_{nm}(t)| \ll |\omega_{nm}|$.
\item Gaussian reservoir ($T_\mu$ field in linear approximation).
\item Slow evolution of the universe: $H_0 \ll \min(|\omega_{nm}|, 1/\tau_S)$.
\end{itemize}

\textbf{Transition between articles}: This master equation establishes the complete mathematical basis for the causal decoherence and objective wave function collapse mechanism discussed in Article I (phenomenology) and connects with the causal genesis derived in Article III through the specific coupling to the $T_\mu$ field with its gradient and precession properties.
\end{derivation}

% --- BLOCK 4: Applications and Physical Consequences of RRT ---

\section*{BLOCK 4: Applications and Physical Consequences of RRT}

\textbf{Transition to Phenomenology}: With the complete mathematical formalism established, we now apply Referential Relativity Theory to specific physical systems, demonstrating its predictive power and observable consequences. The primary novelties are: (1) decoherence rates that grow with redshift, (2) angular dependence in transitions, (3) cosmic precession of preferred directions.

\section{The Consequence of Shielding for Local Quantum Systems}

\begin{derivation}[Nullity of Effects in the Laboratory]
Due to the shielding function $K(\rho)$, we predict that quantum systems in terrestrial laboratories (where the local density $\rho_{lab} \sim 10^{23} \rho_{crit}$) operate in the deep saturation regime.

The effective decoherence rate induced by RRT in the laboratory is exponentially suppressed:
\begin{equation}
\Gamma_{RRT}^{lab} \approx \Gamma_{cosmic} \cdot e^{-\rho_{lab}/\rho_{crit}} \approx 0
\end{equation}

This explains the absence of measurable Lorentz violation in atomic precision experiments and the isotropy observed at CERN, confirming that RRT is an infrared (IR) theory, not an ultraviolet (UV) one.
\end{derivation}

% --- REPLACEMENT OF MICROPHYSICAL BLOCK ---

\section{The Consequence of Shielding for Local Quantum Systems and Matter Stability}

\begin{derivation}[Nullity of Laboratory Effects and Matter Stability]
The direct application of the influence functional derived in Section 22 to quantum systems in terrestrial laboratories reveals the critical importance of the shielding function $K(\rho)$.

For a typical experiment on Earth (density $\rho_{lab} \sim 10^{23} \rho_{crit}$), the effective coupling factor is exponentially suppressed:
\begin{equation}
\xi_{eff} = \xi_T \cdot K(\rho_{lab}) \approx \xi_T \cdot e^{-10^{23}} \approx 0
\end{equation}

This mathematical suppression formalizes the fact that \textbf{Phase 1 (Saturation)} protects local quantum mechanics from cosmological viscosity. Consequently:
\begin{itemize}
    \item \textbf{Zeno Effect and Decoherence:} There is no alteration in qubit decoherence rates or atomic lifetimes in the laboratory.
    \item \textbf{Spectroscopy:} No measurable Lamb shift or Lorentz breaking in atomic clocks (confirming ESA data).
    \item \textbf{Quantum Fields:} The theory reduces strictly to the Standard Model at sub-galactic scales.
\end{itemize}

RRT, therefore, predicts that quantum anomalies can only be observed in experiments conducted in \textbf{deep vacuum} (far from galaxies), where $K(\rho) \to 1$.
\end{derivation}

% -----------------------------------------
\section{Future Research Perspectives in RRT}

\begin{derivation}[Priority Directions for Future Development of Referential Relativity Theory]
\textbf{1. Application to Specific Systems:}
\begin{itemize}
\item \textbf{Cosmology}: Effects of causal fields on inflation and large-scale structure formation.
\item \textbf{Condensed Matter}: Quantum transport in topological materials with gradients.
\item \textbf{Quantum Optics}: Spontaneous emission in anisotropic cavities with precession.
\end{itemize}

\textbf{2. Numerical Implementation:}
\begin{equation}
\rho_S(t+\Delta t) = \mathcal{T} \exp\left[ \int_t^{t+\Delta t} d\tau \, \mathcal{L}(\tau) \right] \rho_S(t)
\tag{44.1}
\end{equation}
Development of efficient algorithms for simulating dissipative dynamics with gradients.

\textbf{3. Crucial Experimental Tests:}
\begin{itemize}
\item Precision measurements of decoherence time $\tau_D(t)$ across different orientations and redshifts.
\item Verification of anisotropy in spontaneous emission via high-resolution spectroscopy.
\item Quantum interference measurements across different spatial directions.
\item Tests of Lorentz invariance violation in controlled quantum systems.
\end{itemize}

\textbf{4. Theoretical Extensions:}
\begin{itemize}
\item \textbf{Quantum Gravity}: Coupling to quantum fluctuations of geometry.
\item \textbf{Quantum Field Theory in Curved Spacetime}: Event horizon effects.
\item \textbf{Non-Equilibrium Many-Body Systems}: Quantum thermalization with gradients.
\item \textbf{Quantum Information}: Channel capacity in anisotropic dissipative media.
\end{itemize}

\textbf{5. Interdisciplinary Connections:}
\begin{itemize}
\item \textbf{Quantum Biology}: Decoherence in biological systems and potential causal roles.
\item \textbf{Materials Science}: Effects of anisotropic dissipation in quantum materials.
\item \textbf{Observational Cosmology}: Testing with CMB, LSS, and supernova data.
\item \textbf{Foundations of Quantum Mechanics}: Interpreting quantum measurement via causal decoherence.
\end{itemize}

\textbf{6. Technical Challenges:}
\begin{itemize}
\item Resolution of the horizon problem in quantum cosmology with causal fields.
\item Formulation of the theory in singular spacetimes (black holes).
\item Extension to non-Abelian gauge theories with Lorentz symmetry breaking.
\item Development of renormalization techniques for cosmic open systems.
\end{itemize}

\textbf{7. Long-Term Perspectives:}
\begin{itemize}
\item Full unification with the Standard Model and quantum gravitation.
\item Formulation of the theory in category and topos language.
\item Development of a grounded quantum theory of causality.
\item Application to resolving the problem of time in quantum gravity.
\end{itemize}

\textbf{8. Expected Scientific Impact:}
\begin{equation}
\Delta \mathcal{K} = \frac{\partial \mathcal{L}}{\partial t} + \{\mathcal{L}, H\} + \mathcal{D}(t)[\mathcal{L}]
\tag{44.2}
\end{equation}
Evolution of scientific knowledge under the dynamics of future RRT research.
\end{derivation}

\section{Conceptual and Philosophical Implications of RRT}

\begin{derivation}[Impact on the Foundations of Theoretical Physics]
\textbf{1. Nature of Time:}
RRT proposes a radical reinterpretation: time is not a background parameter but a dynamic field $T_\mu$ with intrinsic directionality that interacts causally with matter.

\textbf{2. Quantum Measurement Problem:}
Provides a natural physical mechanism for wave function collapse via coupling to the causal field $T_\mu$, resolving the measurement problem without ad hoc collapses.

\textbf{3. Arrows of Time:}
Unification of thermodynamic, cosmological, and quantum arrows of time through a single causal coupling.

\textbf{4. Realism and Locality:}
Reassessment of the concepts of physical reality and local action in the presence of non-local causal fields mediating correlations.

\textbf{5. Emergence of the Classical:}
Rigorous mechanism for the emergence of classical behavior from quantum dynamics, with a rate $\propto (1+z)^2$ explaining why the early universe was more quantum.

\textbf{6. Cosmological Principle:}
Modification of the standard cosmological principle to include fundamental anisotropies and precession, reconciling observations with fundamental theory.

\textbf{7. Math-Physics Relationship:}
New insights into the relationship between mathematical structures (geometry, topology) and physical reality through the concept of referential relativity.
\end{derivation}

\section{Proof of Phase 1: CMS (CERN) Data Analysis}

To validate the saturation regime (Phase 1), we analyzed public data from the CMS detector (CERN Open Data), searching for scalar resonances ("Chrons") in the invariant mass of dimuon events.

\begin{figure}[H]
    \centering
    \includegraphics[width=0.85\textwidth]{grafico39a.png}
    \caption{\textbf{Isotropy Audit at the Z-Peak:} Real CMS data analysis (black dots) shows perfect agreement with the Standard Model (red line). The null residual (blue line) confirms that at the 91 GeV scale, temporal viscosity is zero ($K(\rho) \approx 0$), validating the existence of Phase 1 (Saturated/Shielded).}
    \label{fig:cern_proof}
\end{figure}

The null result does not refute RRT; on the contrary, it calibrates the critical density $\rho_{crit}$, confirming that vacuum viscosity is a low-density emergent phenomenon.

This analysis replaces and corrects previous interpretations that suggested high-energy anomalies, reaffirming that RRT fully respects Lorentz symmetry in the Standard Model regime.

\textbf{Note on Effective Density ($\rho_{eff}$):}
For the calculation of the shielding function $K(\rho)$, density $\rho$ must be interpreted as the total effective energy density ($T_{00}$), encompassing not only baryonic rest mass but also the local gravitational field energy density:
\begin{equation}
\rho_{eff} = \rho_{baryon} + \frac{g^2}{8\pi G c^2}
\end{equation}
In the Solar System and terrestrial orbits (LAGEOS), the contribution of the solar and terrestrial gravitational fields maintains $\rho_{eff} \gg \rho_{crit}$, ensuring that the regime remains in \textbf{Phase 1 (Saturated)}, shielding local experiments against cosmological causal flow.

\section*{Final Conclusions}

\textbf{Referential Relativity Theory (RRT)} represents a deep synthesis between Einstein's geometric vision and Bergson's causal conception, mathematically formalized through the temporal vector field $T_\mu$ and the scalar field $\Phi$. Throughout this paper, we have derived:

\textbf{1. Complete Lagrangian Formulation} unifying geometry, matter, and temporal fields with empirically calibrated couplings.

\textbf{2. Quantum Master Equation} describing the dynamics of open systems coupled to the causal reservoir, incorporating the $(1+z)^2$ gradient, the $\alpha=1.198$ anisotropy, and cosmic precession.

\textbf{3. Rigorous Mathematical Foundation} spanning from covariant action to specific applications in quantum systems, with guaranteed internal consistency.

\textbf{4. Quantitative Predictive Power} with testable predictions for current and future experiments, ranging from laboratory tests to cosmological observations.

\textbf{5. Unified Resolution of Anomalies} including dark matter/energy, Hubble tension, cosmic anisotropies, and black hole corrections.

\vspace{0.3cm} 
Unlike quintessence or modified dark energy models, RRT does not introduce new ad hoc fields, but reveals the vector nature of temporal succession via renormalization flow. This fundamental distinction allows the theory to satisfy Ockham's Razor by unifying dynamical anomalies under a single principle of referential causality.
\vspace{0.3cm}

RRT establishes that time is more than a geometric coordinate – it is an active physical force. The spectroscopic validation via SDSS DR16Q with a \textbf{30.36$\sigma$} mark, reinforced by the $1.69^\circ$ Jackknife stability, suggests that the anisotropy signal can be interpreted as an intrinsic property of the vacuum. The inclusion of the Chern-Simons term identifies the $T_\mu$ field as the Goldstone Boson of causality, unifying referential gravity with the Standard Model of particles.

We invite the scientific community to test these predictions through the Critical Field Tests (CFTs) program outlined in Article I, and to contribute to the development of this theory that promises to transform our fundamental understanding of time, causality, and physical reality.

\section*{Acknowledgments}

I thank the collaborators and reviewers who contributed valuable comments and suggestions. This work was developed independently, supported by open academic discussions and critical analysis of public data.

\section*{Conflict of Interest Statement}

The author declares no conflict of interest regarding the content of this article.

\section*{References}

[1] Cortez, J. C. (2025). \textit{RRT - Volume I: Observational Phenomenology and the Temporal Anisotropy Gradient}.

[2] Bergson, H. (1922). \textit{Durée et Simultanéité: À propos de la théorie d'Einstein}. Paris: Librairie Félix Alcan.

[3] Einstein, A. (1916). \textit{Relativity: The Special and General Theory}. Methuen \& Co Ltd.

[4] Lindblad, G. (1976). On the generators of quantum dynamical semigroups. \textit{Communications in Mathematical Physics}, 48(2), 119-130.

[5] Breuer, H. P., \& Petruccione, F. (2002). \textit{The Theory of Open Quantum Systems}. Oxford University Press.

[6] Zurek, W. H. (2003). Decoherence, einselection, and the quantum origins of the classical. \textit{Reviews of Modern Physics}, 75(3), 715-775.

[7] Brout, D., et al. (2022). The Pantheon+ Analysis: Cosmological Constraints. \textit{The Astrophysical Journal}, 938(2), 110.

[8] Abbott, T. M. C., et al. (2024). The Dark Energy Survey: Cosmology results from the 5-year supernova sample. \textit{arXiv:2401.02929}.

[9] Planck Collaboration. (2018). \textit{Planck 2018 results. VI. Cosmological parameters}. Astronomy \& Astrophysics, 641, A6.

[10] CMB-S4 Collaboration. (2023). CMB-S4 Science Book. \textit{arXiv:2306.06317}.

[11] Feynman, R. P., \& Vernon, F. L. (1963). The theory of a general quantum system interacting with a linear dissipative system. \textit{Annals of Physics}, 24, 118-173.

[12] Gardiner, C. W., \& Zoller, P. (2004). \textit{Quantum Noise: A Handbook of Markovian and Non-Markovian Quantum Stochastic Methods}. Springer.

[13] Schlosshauer, M. (2007). \textit{Decoherence and the Quantum-to-Classical Transition}. Springer.

[14] Caldeira, A. O., \& Leggett, A. J. (1983). Path integral approach to quantum Brownian motion. \textit{Physica A}, 121(3), 587-616.

[15] Keldysh, L. V. (1965). Diagram technique for nonequilibrium processes. \textit{Soviet Physics JETP}, 20(4), 1018-1026.

[16] Cortez, J. C. (2025). Referential Relativity Theory: Code and Data Audit. GitHub Repository. Available at: \url{https://github.com/JeanCCortez/RRT-Audit}.

\appendix

\section{Appendix A: Conventions and Notations}

\begin{itemize}
\item \textbf{Metric}: $(-,+,+,+)$
\item \textbf{Units}: $\hbar=c=1$ when not specified.
\item \textbf{Covariant Derivative}: $\nabla_\mu$
\item \textbf{Trace}: $\Tr$
\item \textbf{Thermal State}: $\rho_T = e^{-\beta H_T}/Z$
\item \textbf{Correlation Functions}: $\mathcal{C}(\tau, \tau') = \langle B(\tau)B(\tau') \rangle$
\item \textbf{Spectral Density}: $J(\omega, \theta, t)$
\item \textbf{Redshift}: $z(t)$ related to the scale factor $a(t)$ by $1+z = a_0/a(t)$
\end{itemize}

\section{Appendix B: Computational Code for Simulations}

Available at: \url{https://github.com/JeanCCortez/RRT-Audit}

Includes Python implementations of:
\begin{itemize}
\item Master equation solver with time-dependent coefficients.
\item Calculation of correlation functions with gradients.
\item Galactic rotation curve simulations with variable viscosity.
\item Parameter fitting to observational data.
\end{itemize}

\end{document}
